%Planning to understand/simplify/improve. If you are interested, please contact me (Evangelos Kipouridis), I would be more than happy to do that collaboratively.
%Licensed under CC BY 4.0
\documentclass[11pt]{article}
\usepackage[T1]{fontenc}
\usepackage[utf8]{inputenc}
\usepackage{lmodern,microtype}
\usepackage[a4paper,margin=28mm]{geometry}
\usepackage{amsmath,amssymb,amsthm,mathtools}
\usepackage{enumitem,listings,xcolor}
\usepackage[colorlinks=true,linkcolor=blue!55!black,citecolor=blue!55!black,urlcolor=blue!55!black]{hyperref}
\numberwithin{equation}{section}
\newtheorem{theorem}{Theorem}[section]
\newtheorem{lemma}[theorem]{Lemma}
\theoremstyle{definition}
\newtheorem{definition}[theorem]{Definition}
\newcommand{\Z}{\mathbb Z}
\newcommand{\Q}{\mathbb Q}
\newcommand{\R}{\mathbb R}
\newcommand{\E}{\mathbb E}
\newcommand{\Prob}{\mathbb P}
\newcommand{\Var}{\operatorname{Var}}
\newcommand{\ind}{\mathbf 1}
\newcommand{\ecc}{\operatorname{ecc}}
\newcommand{\rad}{\operatorname{rad}}
\newcommand{\cL}{\overline L}
\newcommand{\pos}[1]{\left(#1\right)_{+}}
\DeclareMathOperator*{\argmin}{arg\,min}
\lstset{basicstyle=\ttfamily\scriptsize,columns=fullflexible,
 keepspaces=true,showstringspaces=false,breaklines=true,
 breakatwhitespace=false,frame=single,rulecolor=\color{black!25},
 numbers=left,numberstyle=\tiny\color{black!50},numbersep=6pt,
 xleftmargin=10pt,framexleftmargin=8pt,tabsize=4}
\setlist{itemsep=3pt,topsep=5pt}
\title{Upper and Lower Bounds for Approximating Graph Burning}
\author{}
\date{}

\begin{filecontents*}[overwrite]{graph_burning_profile.json}
{
  "a": "0.11214",
  "atom": "0.55101",
  "knots": [
    "0",
    "0.21324",
    "0.64319",
    "1.11214"
  ],
  "density": [
    "1.35164",
    "0.94044",
    "1.35164"
  ]
}
\end{filecontents*}

\begin{filecontents*}[overwrite]{graph_burning_certificate.py}
#!/usr/bin/env python3
"""Exact scalar certificate for a rational piecewise constant enlargement density."""
from fractions import Fraction as F
from math import factorial,lcm
from array import array
from pathlib import Path
from decimal import Decimal,localcontext
import argparse,json,subprocess,sys,time

S=10**12;T=10**40
def floorf(x):return x.numerator//x.denominator
def ceilf(x):return -((-x.numerator)//x.denominator)
def cd(a,b):return -((-a)//b)
ECU=[ceilf(F((-1)**j,factorial(j))*T) for j in range(21)]
ECL=[floorf(F((-1)**j,factorial(j))*T) for j in range(22)]
FCU=[ceilf(F((-1)**j,factorial(j+1))*T) for j in range(21)]
FCL=[floorf(F((-1)**j,factorial(j+1))*T) for j in range(22)]
def evaluate(num,den,coef,upper):
    x=num*T//den if upper else cd(num*T,den);y=coef[-1]
    for z in reversed(coef[:-1]):y=(cd(y*x,T) if upper else y*x//T)+z
    return y
def up_exp(num,den):return min(S,cd(evaluate(num,den,ECU,True)*S,T))
def lo_exp(num,den):return max(0,evaluate(num,den,ECL,False)*S//T)
def up_phi(num,den):return min(S,cd(evaluate(num,den,FCU,True)*S,T))

def main():
    ap=argparse.ArgumentParser(description=__doc__)
    ap.add_argument('profile')
    ap.add_argument('--source-grid',type=int,default=65536)
    ap.add_argument('--time-grid',type=int,default=65536)
    ap.add_argument('--primitive-grid',type=int,default=524288)
    ap.add_argument('--reuse-tables',action='store_true')
    args=ap.parse_args();nr=args.source_grid;nt=args.time_grid;nq=args.primitive_grid
    assert nq%nr==nq%nt==nt%nr==0
    path=Path(args.profile).resolve();base=path.parent
    obj=json.loads(path.read_text());a=F(obj['a']);atom=F(obj['atom'])
    knots=list(map(F,obj['knots']));dens=list(map(F,obj['density']));n=len(dens)
    assert knots[0]==0 and knots[-1]==1+a and len(knots)==n+1
    assert all(x<y for x,y in zip(knots,knots[1:])) and all(x>=0 for x in dens)
    W=atom+sum(p*(u-l) for p,l,u in zip(dens,knots,knots[1:]))
    assert 0<a<F(3,20) and 0<atom<W/2 and W<2 and max(dens)<2
    assert atom+max(dens)<2 and atom+max(dens)+2*max(dens)*W<7
    assert max(dens)*a<1
    coord=lcm(*(x.denominator for x in knots+[a]));kc=[int(x*coord) for x in knots]
    atomS=int(atom*S);WS=int(W*S);PS=[int(p*S) for p in dens]
    assert atom*S==atomS and W*S==WS and all(p*S==q for p,q in zip(dens,PS))
    assert all(p%coord==0 for p in PS)
    assert 2*WS*nq*(coord*nt)**2 < 2**126
    FK=[atomS]
    for p,l,u in zip(PS,kc,kc[1:]):
        assert p*(u-l)%coord==0
        FK.append(FK[-1]+p*(u-l)//coord)
    assert FK[-1]==WS
    def fnum(x,nn):
        z=atomS*coord*nn
        for p,l,u in zip(PS,kc,kc[1:]):z+=p*min(max(x-l*nn,0),(u-l)*nn)
        return z
    def gnum(j,nn):
        if j*coord>=kc[-1]*nn:return 0
        mx=2*fnum(j*coord,nn)
        for k,f in zip(kc,FK):
            if k*nn<=2*j*coord:mx=max(mx,f*coord*nn+fnum(2*j*coord-k*nn,nn))
        return max(0,min(WS*coord*nn,2*WS*coord*nn-mx))
    def dnum(j,nn):return WS*coord*nn-fnum(j*coord,nn)
    stem=path.stem+f'_{nr}_{nt}_{nq}';tables=base/(stem+'.bin');result=base/(stem+'.certificate.json')
    if not(args.reuse_tables and tables.exists()):
        started=time.monotonic()
        qu=array('q');ql=array('q');fu=array('q');den=S*coord*nq
        for j in range(3*nq+1):
            z=gnum(j,nq)
            qhi=qlo=fhi=S if z==0 else 0
            if z:qhi=up_exp(z,den);qlo=lo_exp(z,den);fhi=up_phi(z,den)
            assert 0<=qlo<=qhi<=S and 0<fhi<=S
            qu.append(qhi);ql.append(qlo);fu.append(fhi)
        assert all(x<=y for x,y in zip(qu,qu[1:])) and all(x<=y for x,y in zip(ql,ql[1:]))
        assert all(x<=y for x,y in zip(fu,fu[1:]))
        ou=array('q',[up_exp(dnum(j,nt),S*coord*nt) for j in range(3*nt+1)])
        eu=array('q',[up_exp(gnum(j,2*nt),S*coord*2*nt) for j in range(2*nt+1)])
        assert sys.byteorder=='little' and array('q').itemsize==8
        with tables.open('wb') as out:
            array('q',[211250002,S,nq,nt,nr,coord,n,atomS,WS,int(a*coord),int((1+a)*coord)]).tofile(out)
            for vals in (array('q',kc),array('q',PS),qu,ql,fu,ou,eu):vals.tofile(out)
        print('Exact tables generated in',round(time.monotonic()-started,2),'seconds',flush=True)
    executable=base/'graph_burning_certificate_verify'
    subprocess.run(['g++','-std=c++17','-O3','-Wall','-Wextra',str(base/'graph_burning_certificate.cpp'),'-o',str(executable)],check=True)
    subprocess.run([str(executable),str(tables),str(result)],check=True)
    out=json.loads(result.read_text());j=F(int(out['numerator']),int(out['denominator']));cap=F(int(out['capacity_numerator']),int(out['capacity_denominator']))
    assert j<F(999999,1000000) and cap>F(1,10)
    print('J_UPPER_EXACT =',j)
    with localcontext() as ctx:
        ctx.prec=60
        print('J_UPPER_DECIMAL =',Decimal(j.numerator)/Decimal(j.denominator))
        print('A_LOWER_DECIMAL =',Decimal(cap.numerator)/Decimal(cap.denominator))
    print('CERTIFIED: J < 0.999999 and A(s) > 0.1.')
if __name__=='__main__':main()
\end{filecontents*}

\begin{filecontents*}[overwrite]{graph_burning_certificate.cpp}
#include <algorithm>
#include <cassert>
#include <cstdint>
#include <fstream>
#include <iostream>
#include <string>
#include <vector>
using i64=std::int64_t;using i128=__int128_t;
static i64 cd(i128 a,i128 b){assert(b>0);i128 v=a>=0?(a+b-1)/b:a/b;assert(v<=INT64_MAX&&v>=INT64_MIN);return (i64)v;}
static std::string dec(i128 x){if(!x)return "0";bool neg=x<0;if(neg)x=-x;std::string s;while(x){s.push_back('0'+x%10);x/=10;}if(neg)s.push_back('-');std::reverse(s.begin(),s.end());return s;}
static i64 ri(std::ifstream&f){i64 v;f.read((char*)&v,8);assert(f.good());return v;}
static std::vector<i64> rv(std::ifstream&f,size_t n){std::vector<i64>v(n);f.read((char*)v.data(),8*n);assert(f.good());return v;}
int main(int argc,char**argv){
 assert(argc==3);std::ifstream f(argv[1],std::ios::binary);assert(f.good());
 assert(ri(f)==211250002);const i64 S=ri(f),nq=ri(f),nt=ri(f),nr=ri(f),C=ri(f),n=ri(f);
 const i64 atom=ri(f),W=ri(f),ac=ri(f),mc=ri(f);
 assert(S==1000000000000LL&&nq%nt==0&&nq%nr==0&&nt%nr==0&&n>0&&n<30);
 auto knots=rv(f,n+1),P=rv(f,n),qu=rv(f,3*nq+1),ql=rv(f,3*nq+1),fu=rv(f,3*nq+1),ou=rv(f,3*nt+1),enu=rv(f,2*nt+1);
 const i64 H=S*nq,T=C*nq,AH=atom*nq,WH=W*nq;
 assert(H>0&&H<1000000000000000000LL&&T>0&&T<1000000000000000000LL&&T%nt==0&&T%nr==0);
 std::vector<i64> slope(n),pu(3*nq+1),pl(3*nq+1);
 for(int h=0;h<n;h++){assert(P[h]%C==0);slope[h]=P[h]/C;}
 for(i64 j=1;j<=3*nq;j++){
  assert(0<=ql[j]&&ql[j]<=qu[j]&&qu[j]<=S&&qu[j-1]<=qu[j]&&ql[j-1]<=ql[j]&&fu[j-1]<=fu[j]&&fu[j]<=S);
  pu[j]=pu[j-1]+qu[j];pl[j]=pl[j-1]+ql[j-1];
 }
 auto pf=[&](const std::vector<i64>&pref,const std::vector<i64>&qs,i64 x,bool right,bool up){
  assert(x>=0&&x<3*T);i64 j=x/C,z=x%C;
  return pref[j]+(up?cd((i128)qs[j+(right?1:0)]*z,C):(i64)((i128)qs[j+(right?1:0)]*z/C));
 };
 const i64 nj=cd((i128)mc*nt,C),step=T/nt;
 std::vector<i128> acc(nj+1);
 for(i64 ii=1;ii<=nr;ii++){
  const i64 rc=ii*(T/nr),rq=ii*(nq/nr);
  const i64 aqlo=(i64)((i128)AH*ql[rq]/S),aqup=cd((i128)AH*qu[rq],S);
  std::vector<i64>lo(n),hi(n),pustart(n),plstart(n),plend(n);
  for(int h=0;h<n;h++){
   lo[h]=rc+knots[h]*nq;hi[h]=rc+knots[h+1]*nq;
   pustart[h]=pf(pu,qu,lo[h],true,false);plstart[h]=pf(pl,ql,lo[h],false,true);plend[h]=pf(pl,ql,hi[h],false,false);
  }
  auto tail=[&](i64 x){i64 v=WH-(x>rc?AH:0);for(int h=0;h<n;h++)v-=slope[h]*std::clamp(x-lo[h],(i64)0,hi[h]-lo[h]);assert(v>=0&&v<=WH);return v;};
  auto preup=[&](i64 x){i64 v=x>rc?aqup:0;for(int h=0;h<n;h++)if(x>lo[h])v+=cd((i128)P[h]*(pf(pu,qu,std::min(x,hi[h]),true,true)-pustart[h]),S);return v;};
  auto prelo=[&](i64 x){i64 v=x>rc?aqlo:0;for(int h=0;h<n;h++)if(x>lo[h])v+=(i64)((i128)P[h]*std::max<i64>((x>=hi[h]?plend[h]:pf(pl,ql,x,false,false))-plstart[h],0)/S);return v;};
  i64 running=preup(mc*nq)+tail(mc*nq);
  for(i64 j=nj;j>=0;j--){
   const i64 x=j*step,qq=j*(nq/nt),ta=tail(x);
   const i64 point=preup(x)+cd((i128)fu[qq]*ta,S);
   running=std::max(running,point);
   // In an open cell beginning at the source atom, the atom has already
   // moved from T into Q. Treat the endpoint itself by `point` above.
   const i64 tright=ta-(x==rc?AH:0);
   const i64 cell=preup(x+step)+cd((i128)fu[qq+nq/nt]*tright,S);
   running=std::max(running,cell);
   i64 val=std::min(ta,std::max<i64>(0,running-prelo(x)));
   if(x>=mc*nq)val=ta;
   acc[j]+=val;
  }
  if(ii%(std::max<i64>(1,nr/8))==0)std::cerr<<"source rows "<<ii<<" / "<<nr<<"\n";
 }
 const i64 TC=C*nt;const i128 denmass=2*(i128)TC*TC;
 auto cdf=[&](i64 x){
  auto pp=[](i64 z)->i128{return z>0?(i128)z*z:0;};
  i128 val=2*(i128)AH*TC*std::clamp(x,(i64)0,TC);
  for(int h=0;h<n;h++){
   const i64 l=knots[h]*nt,u=knots[h+1]*nt;
   val+=(i128)P[h]*nq*(pp(x-l)-pp(x-u)-pp(x-TC-l)+pp(x-TC-u));
  }
  return val;
 };
 std::vector<i64> old(3*nt+1);i128 nxt=cdf(3*TC);
 for(i64 j=3*nt-1;j>=0;j--){i128 cur=cdf(j*C);assert(nxt>=cur);i64 mass=cd(nxt-cur,denmass);old[j]=old[j+1]+cd((i128)mass*ou[j+1],S);nxt=cur;}
 assert(nxt==0&&cdf(3*TC)==(i128)WH*denmass);
 assert((i128)ac*H%C==0);const i64 abase=2*H+(i64)((i128)ac*H/C),tstep=H/nt;
 i64 running=INT64_MAX,amin=INT64_MAX,previous_raw=0;i128 jsum=0;
 for(i64 j=0;j<=2*nt;j++){
  i64 mu=old[j];if(j<=nj)mu=std::min(mu,cd(acc[j],nr));
  const i64 raw=abase-j*tstep-mu;
  running=std::min(running,raw);
  if(j){
   // Both mean envelopes lose at most Lambda([x,y]) when their threshold
   // increases from x to y. If Lambda's density is at most one on a cell,
   // 2+a-theta-M(theta) is nonincreasing there: its minimum is the right
   // endpoint. Otherwise use the left endpoint minus one cell width.
   const i64 left=(j-1)*C,right=j*C;
   i128 density_upper=left<TC?(i128)atom*TC:0;
   for(int h=0;h<n;h++){
    const i64 l=knots[h]*nt,u=knots[h+1]*nt;
    const i64 overlap=std::max<i64>(0,std::min(TC,right-l)-std::max<i64>(0,left-u));
    density_upper+=(i128)P[h]*overlap;
   }
   const i64 cell_lower=density_upper<=(i128)S*TC?raw:previous_raw-tstep;
   running=std::min(running,cell_lower);
  }
  previous_raw=raw;
  const i64 cap=running;amin=std::min(amin,cap);assert(cap>H/10);
  if(j)jsum+=cd((i128)enu[j]*H,cap);
 }
 const i128 denJ=2*(i128)nt*S;
 std::cerr<<"J exact: "<<dec(jsum)<<" / "<<dec(denJ)<<"\n";
 std::ofstream out(argv[2]);out<<"{\"numerator\":\""<<dec(jsum)<<"\",\"denominator\":\""<<dec(denJ)<<"\",\"capacity_numerator\":\""<<amin<<"\",\"capacity_denominator\":\""<<H<<"\"}\n";
 assert(jsum*1000000<999999*denJ);
 std::cerr<<"CERTIFIED J < 999999/1000000 and A > 1/10\n";
}
\end{filecontents*}

\begin{document}
\maketitle
\begin{abstract}
Let $G=(V,E)$ be a finite nonempty simple undirected graph, let
$n=|V|$, and let $B_G(v,r)$ denote the vertices at graph distance at most
$r$ from $v$. The burning number $b(G)$ is the least positive integer
$L$ for which $V=\bigcup_{i=1}^{L}B_G(v_i,L-i)$ for some vertices
$v_1,\ldots,v_L$. We give a deterministic polynomial-time
$(2.11214+\varepsilon)$-approximation for every fixed $\varepsilon>0$;
in particular, a legal burning sequence of length at most $\lfloor2.112141b(G)\rfloor$
can be constructed. The numerical inequality used in this proof is
certified by exact rational arithmetic, with the complete verification
procedure included. Given a positive integer $k$, we also give a $4^k n^{O(1)}$-time
algorithm returning a legal burning sequence of length at most $2k-1$
whenever $b(G)\le k$.
Finally, for every fixed $\alpha<2$ with $\alpha\ge1$, approximation
within factor $\alpha$ is NP-hard even on connected bipartite graphs.
The hardness proof is a reduction from the constant-gap Set Cover problem.
\end{abstract}

\section{Definitions and theorem statements}\label{sec:definitions}

All graphs in this paper are finite, simple, undirected, and nonempty.
Thus a graph $G=(V,E)$ consists of a nonempty vertex set $V$ and a set
$E$ of unordered pairs of distinct vertices. Every edge has length one.
Fix an ordering of $V$ and the induced lexicographic orderings of finite
tuples. Whenever a choice is not unique, use the first permitted
candidate in the corresponding fixed ordering.
Write $n=|V|$. The distance $d_G(u,v)$ is the minimum number of edges in
a path from $u$ to $v$, and is $+\infty$ when no such path exists.
A connected component is an equivalence class of the relation
$d_G(u,v)<+\infty$. A graph is connected if it has one component.
A graph is bipartite if its vertices admit a partition into two sets
such that every edge joins different sets. For $v\in V$ and $r\ge0$, set
\[
 B(v,r)=B_G(v,r)=\{w\in V:d_G(v,w)\le r\}.
\]
For a component $C$ and $v\in C$, define
\[
 \ecc_C(v)=\max_{w\in C}d_G(v,w),\qquad
 \rad(C)=\min_{v\in C}\ecc_C(v).
\]

For an integer $q\ge0$, write $[q]=\{1,\ldots,q\}$, with $[0]=\varnothing$.
The symbols $\Z$, $\Q$, and $\R$ denote the integers, rational numbers,
and real numbers, respectively. A subscript restricts these sets by the
displayed inequality; for example, $\Z_{\ge1}$ denotes the positive
integers.
For a proposition $P$, its indicator $\ind_P$ is one if $P$ holds and zero
otherwise. Put $x_+=\max\{0,x\}$ and
$[x]_0^c=\max\{0,\min\{c,x\}\}$ for $c\ge0$.
The symbols $\lfloor x\rfloor$ and $\lceil x\rceil$ denote the greatest
integer at most $x$ and the least integer at least $x$, respectively.
All finite decimal constants denote exact rational numbers, and $\log$
denotes the natural logarithm. For a function on a finite nonempty set,
$\argmin$ denotes the set of arguments attaining its minimum.

A finite probability space is a finite set $\Omega$ with nonnegative
weights $\pi(\omega)$ summing to one. For $A\subseteq\Omega$ and
$Z:\Omega\to\R$, write
$\Prob(A)=\sum_{\omega\in A}\pi(\omega)$,
$\E Z=\sum_{\omega\in\Omega}\pi(\omega)Z(\omega)$, and
$\Var Z=\E(Z-\E Z)^2$.
Conditional expectation on an event of positive probability is the
weighted average restricted to that event and divided by its probability.
Conditional expectation given specified random variables applies this
rule separately for each joint value of those variables having positive
probability.
Random variables are independent when their joint point probabilities
are products of their individual point probabilities.

\begin{definition}\label{def:burning}
A covering sequence of length $L\in\Z_{\ge1}$ is a sequence
$(v_1,\ldots,v_L)\in V^L$ such that
\begin{equation}\label{eq:cover-sequence}
 V=\bigcup_{i=1}^{L}B(v_i,L-i).
\end{equation}
The burning number $b(G)$ is the minimum length of a covering sequence.
The vertex chosen at a round is called the source at that round.
For a sequence $(v_1,\ldots,v_L)$, define its burned sets by
\[
 \mathcal B_0=\varnothing,\qquad
 \mathcal B_i=\bigcup_{j=1}^i B(v_j,i-j)\quad(i\in[L]).
\]
The set $\mathcal B_i$ consists of the vertices burned by the end of
round $i$. A legal burning sequence is a covering sequence satisfying
$v_i\notin\mathcal B_{i-1}$ for every $i\in[L]$.
This condition implies $\mathcal B_i\ne V$ for $i<L$, so the sequence
terminates at the first round when all vertices have burned.
\end{definition}

Lemma~\ref{lem:legalization} proves that the same minimum is obtained
from legal burning sequences. In particular, $1\le b(G)\le n$.
An $\alpha$-approximation, for $\alpha\ge1$, returns a legal burning
sequence of length at most $\alpha b(G)$ on every input graph.
Running time refers to bit operations on an adjacency-matrix encoding of
$G$ and binary encodings of integer or rational parameters. The notation
$n^{O_\varepsilon(1)}$ permits both a multiplicative constant and an
exponent depending on a fixed $\varepsilon$, but not on $G$ or $b(G)$.
The notation $4^k n^{O(1)}$ has an exponent independent of $k$.

\begin{theorem}\label{thm:poly}
For every fixed $\varepsilon>0$, there is a deterministic
$n^{O_\varepsilon(1)}$-time $(2.11214+\varepsilon)$-approximation for
graph burning. In particular, there is a deterministic polynomial-time
algorithm returning a legal burning sequence of length at most
$\lfloor2.112141b(G)\rfloor$.
\end{theorem}

\begin{theorem}\label{thm:fpt}
Given $G$ and $k\in[n]$, a deterministic algorithm runs in
$4^k n^{O(1)}$ time and either returns a legal burning sequence of length
at most $2k-1$ or returns \textnormal{NO}. If it returns \textnormal{NO},
then $b(G)>k$. Consequently there is a deterministic
$4^{b(G)}n^{O(1)}$-time algorithm returning a legal burning sequence of
length at most $2b(G)-1$.
\end{theorem}

We use the conventional complexity classes $\mathrm P$ and $\mathrm{NP}$:
$\mathrm P$ consists of decision problems solvable in deterministic
polynomial time, and $\mathrm{NP}$ consists of decision problems whose
positive answers have polynomial-size certificates verifiable in
polynomial time. A promise problem specifies disjoint sets of positive
and negative inputs, and an algorithm must be correct on those inputs.
NP-hardness of a promise problem means that every problem in
$\mathrm{NP}$ has a polynomial-time many-one reduction mapping positive
inputs to positive inputs and negative inputs to negative inputs.

\begin{theorem}\label{thm:hardness}
For every fixed integer $\ell\ge3$, the following promise problem is
NP-hard on connected bipartite graphs. Its input is a graph $G$ and a
positive integer $k$, and its two cases are
\begin{equation}\label{eq:hardness-gap}
 b(G)\le 2(\ell+1)k
 \qquad\text{and}\qquad
 b(G)>(4\ell-2)k.
\end{equation}
Therefore, for every fixed $\alpha\in[1,2)$, a polynomial-time
$\alpha$-approximation on connected bipartite graphs implies
$\mathrm P=\mathrm{NP}$.
\end{theorem}

Only these three results and the lemmas needed to prove them are developed
below. The polynomial-time upper bound uses geometric grouping of radii
and cumulative counts; this organization is related to the grouping of
Blauth, N\"obel, and Zenklusen~\cite{BNZ}. All rounding statements used
here are proved here.

\section{Covering sequences, deadlines, and cumulative counts}
\label{sec:tools}

\begin{lemma}[Legalization]\label{lem:legalization}
Let $m,L$ be positive integers. Suppose the balls $B(u_j,R_j)$,
$j\in[m]$, cover $V$ and have pairwise distinct assigned rounds
$t_j\in[L]$, where $0\le R_j\le L-t_j$. A legal burning sequence of
length at most $L$ can be found in polynomial time.
\end{lemma}
\begin{proof}
Fill every unassigned round with any vertex to obtain a planned sequence
$(v_1,\ldots,v_L)$ satisfying~\eqref{eq:cover-sequence}. Process its
rounds in increasing order. If every vertex has burned, terminate. If
$v_i$ has not burned, ignite $v_i$. Otherwise ignite the first vertex
that has not burned. In the latter case, some actual source $w$ ignited
at an earlier round $j<i$ satisfies $d_G(w,v_i)\le i-j$.
The triangle inequality gives
\[
 B(v_i,L-i)\subseteq B(w,L-j).
\]
If the procedure reaches round $L$, every planned ball is therefore
contained in the radius-$(L-j)$ ball of some actual source at round $j$.
The planned balls cover $V$, so all vertices have burned by round $L$.
Termination at an earlier round also gives the required sequence.
Padding is represented implicitly. Each processed round ignites an
unburned vertex, so at most $n$ rounds are processed. Distances and
all choices can be computed in polynomial time.
\end{proof}

\subsection{The parameterized algorithm}

\begin{proof}[Proof of Theorem~\ref{thm:fpt}]
Partition $[k]$ into consecutive blocks as follows. For a current prefix
$[q]$, take its last block
$[h,q]$, where $h=\lfloor q/2\rfloor+1$, and repeat on $[h-1]$ until
the prefix is empty. The notation $[h,q]$ denotes
$\{h,h+1,\ldots,q\}$. Replace each radius $k-i$ in a block $[h,q]$ by
$k-h$. These are rounded radii, each occurring with the multiplicity of
its block.

Maintain uncovered vertices and the multiset of unused rounded radii.
Choose the first uncovered vertex $v$. For each distinct unused rounded
radius $r$, form a branch that chooses $B(v,2r)$, deletes its vertices
from the uncovered set, and deletes one occurrence of $r$. A branch
succeeds when no vertex remains. It fails when vertices remain but no
radius remains. Return \textnormal{NO} if every branch fails.

If $b(G)\le k$, there is a cover by the original radii $k-i$, and hence
by the rounded radii. In a branch consistent with such a cover, let an
unused covering ball $B(w,r)$ contain the chosen vertex $v$. Then
$B(w,r)\subseteq B(v,2r)$. Removing this covering ball and the vertices
of $B(v,2r)$ leaves a cover of the remaining vertices by the other unused
covering balls. Induction on the remaining radii proves that some branch
succeeds. Thus \textnormal{NO} implies $b(G)>k$.

Assign balls produced from a block $[h,q]$ to distinct rounds in that
block. Since $q\le2h-1$,
\[
 (2k-1)-q\ge2(k-h).
\]
Every chosen ball fits its assigned round in horizon $2k-1$.
Lemma~\ref{lem:legalization} gives the legal sequence.

For each prefix size $q$ occurring in the recursive partition, let $M(q)$
be the number of permutations of the rounded-radius multiset from the
first $q$ indices. The blocks have distinct rounded radii, and
\[
 M(q)=\binom{q}{\lceil q/2\rceil}M(\lfloor q/2\rfloor),
 \qquad M(0)=1.
\]
The bound $\binom{q}{\lceil q/2\rceil}\le2^q$ and
$\sum_{j\ge0}\lfloor k/2^j\rfloor\le2k$ imply $M(k)\le4^k$.
Every search node is a prefix of a multiset permutation. The number of
nodes is at most $(k+1)M(k)$, and each node requires polynomial time.
This proves the first assertion. Test $k=1,\ldots,n$ and stop at the
first success. Success occurs by $b(G)$; the sum of running times through
that guess is $4^{b(G)}n^{O(1)}$, and the returned length is at most
$2b(G)-1$.
\end{proof}

\subsection{Unit jobs with deadlines}

A unit job occupies one integer position and has an integer deadline
$d$; its position must be at most $d$. A placement of a family of unit
jobs assigns pairwise distinct positions to those jobs.

\begin{lemma}\label{lem:deadlines}
Let $T$ be a positive integer. Consider a finite family of unit jobs
with deadlines in $[T]$, and let $m_j$ be the number
with deadline at most $j$ for $j\in[T]$, with $m_0=0$.
They can be placed in $[T]$ if and only if $m_j\le j$ for every $j\in[T]$.
Placing each job in the latest
unoccupied position at most its deadline produces a placement for which
the number $F(D)$ of free positions in $[D]$ satisfies, for every
integer $D\in\{0,\ldots,T\}$,
\begin{equation}\label{eq:free-positions}
 F(D)=\min_{D\le j\le T}(j-m_j)\qquad(0\le D\le T).
\end{equation}
\end{lemma}
\begin{proof}
Necessity of $m_j\le j$ follows by counting positions. For sufficiency,
sort the deadlines as $d_1\le\cdots\le d_s$. Then $i\le m_{d_i}\le d_i$,
so placing job $i$ at position $i$ is feasible.
The latest-position procedure also succeeds: starting with any feasible
placement of the unplaced jobs, move the next job from its current
position to the latest free position allowed by its deadline. If another
unplaced job occupies that position in the feasible placement, move that
job to the first job's old position. The latter position is earlier, so
both deadlines are respected. This preserves a feasible completion.

For every feasible placement and $j\ge D$, at least
$m_j-(j-D)$ jobs occupy $[D]$, so $F(D)\le j-m_j$.
For the reverse inequality, let $u$ be the first free position after
$D$ in the latest-position placement, if one exists. Every job placed
before $u$ has deadline less than $u$, since $u$ was available when that
job was placed. Thus $F(D)=u-1-m_{u-1}$, and $u-1\ge D$.
If no free position follows $D$, then $F(D)=T-m_T$.
In either case equality holds for one of the terms in the minimum.
\end{proof}

\subsection{Fractional covering with cumulative counts}

\begin{lemma}\label{lem:cumulative-rounding}
Let $U\subseteq V$ be a client set, and let $s$ be a positive integer.
Let $R_1>\cdots>R_s\ge0$ be
integer radii and $1\le a_1\le\cdots\le a_s$ be integers.
For a finite index set $\mathcal E$, let $c_e\in V$, let
$\iota(e)\in[s]$, and let $z_e\in\Q_{\ge0}$ for each $e\in\mathcal E$.
Suppose
\begin{equation}\label{eq:rounding-hypotheses}
 \sum_{e:w\in B(c_e,R_{\iota(e)})}z_e\ge1\quad(w\in U),
 \qquad
 \sum_e\frac{z_e}{a_{\iota(e)}}\le1.
\end{equation}
There is a deterministic polynomial-time cover of $U$ by balls
$B(u,2R_j)$, with at most $a_i$ chosen balls having type $j\le i$,
simultaneously for every $i\in[s]$.
\end{lemma}
\begin{proof}
Merge consecutive types having the same capacity into their largest
radius, keeping each center and mass. This preserves the weighted sum
and can only increase coverage. We may assume that the capacities
strictly increase. Put $a_0=0$.

Maintain the uncovered client set, still denoted by $U$, and delete
fractional balls with empty intersection with $U$. In phase $j$, choose
at most $a_j-a_{j-1}$ centers. Choose $u\in U$ maximizing
\begin{equation}\label{eq:maximum-release}
 \sum_{\substack{e\text{ remaining}\\\iota(e)\ge j}}
 \frac{z_e}{a_{\iota(e)}}
 \ind_{B(c_e,R_{\iota(e)})\cap U\subseteq B(u,2R_j)}.
\end{equation}
Choose $B(u,2R_j)$, remove its clients from $U$, and delete all fractional
balls counted in~\eqref{eq:maximum-release}. At the end of phase $j$,
discard type $j$. Stop if $U$ becomes empty. The displayed weights are
kept fixed; only the remaining balls and clients change.

We prove coverage by induction on $s$. In the first phase, define
\[
 \Phi(w)=\sum_{e:\iota(e)>1,\ w\in B(c_e,R_{\iota(e)})}
             (1-a_1/a_{\iota(e)})z_e.
\]
Every remaining ball containing $w$ lies in $B(w,2R_1)$, and its
weighted mass is at least $(1-\Phi(w))/a_1$ by
\eqref{eq:rounding-hypotheses}. Therefore a maximum in
\eqref{eq:maximum-release} deletes at least
$(1-\min_{w\in U}\Phi(w))/a_1$ weighted mass.
Deleted fractional balls contain no surviving client. Thus $\Phi(w)$
does not change for a client that survives a deletion.

Suppose clients remain after $a_1$ choices and let $\zeta$ be their
minimum value of $\Phi$. At each choice the then-current minimum was
at most $\zeta$. Remaining weighted mass is at most
$1-a_1(1-\zeta)/a_1=\zeta$. If $\zeta=0$, this contradicts retained
fractional coverage. Otherwise discard type one and set
\[
 z'_e=\frac{a_{\iota(e)}-a_1}{a_{\iota(e)}\zeta}z_e,
 \qquad a'_i=a_i-a_1\quad(i>1).
\]
New coverage is $\Phi(w)/\zeta\ge1$, and new weighted mass is at most
one. The induction hypothesis applies. Moreover
$z'_e/a'_{\iota(e)}=(z_e/a_{\iota(e)})/\zeta$, so its
maximum-release choices agree with the stated algorithm. The phase
quotas give every cumulative bound.
Each choice removes its chosen client, so at most $|U|$ choices are
made. All tests and rational sums are polynomial-time computations.
\end{proof}

The rule also counts fractional balls that do not contain the chosen
center. In particular, if $B(v,r)$ meets $B(u,R)$ and $r\le R/2$, then
$d_G(u,v)\le R+r$ and $B(v,r)\subseteq B(u,2R)$.
Thus any number of such balls can be deleted in one radius-$2R$ choice.
More generally, let $m\ge0$ be an integer, let $R,s\ge0$, let $u\in V$, and let
$\{B(v_j,r_j):j\in[m]\}$ be a finite family with $0\le r_j\le s$
and $B(v_j,r_j)\cap B(u,R)\ne\varnothing$ for every $j$.
For each $j$, the triangle inequality gives
$B(v_j,r_j)\subseteq B(u,R+2s)$.
Consequently one ball of radius $R+2s$ covers $B(u,R)$ and all balls
in this family, regardless of $m$.

\section{A fractional cover and independent proposals}\label{sec:proposals}

Fix an integer guess $k\in[n]$. The upper-bound construction uses the
following rational constants:
\begin{equation}\label{eq:profile}
\begin{gathered}
 a=\frac{5607}{50000},\qquad A_0=\frac{55101}{100000},\qquad
 c_0=0,\quad c_1=\frac{5331}{25000},\quad
 c_2=\frac{64319}{100000},\quad c_3=1+a,\\
 p_1=p_3=\frac{33791}{25000},\qquad
 p_2=\frac{23511}{25000},\qquad p_{\max}=p_1.
\end{gathered}
\end{equation}
Define a finite measure $\mu$ on $\R_{\ge0}$ by the integration rule
\begin{equation}\label{eq:measure}
 \int f(h)\,d\mu(h)
 = A_0 f(0)+\sum_{j=1}^3p_j\int_{c_{j-1}}^{c_j}f(h)\,dh.
\end{equation}
The rule is used for the bounded piecewise continuous functions appearing
below. A positive mass assigned to a single point is called an atom at
that point. The measure $\mu$ has an atom of mass $A_0$ at zero, and
no other point has positive mass.
Its total mass, its cumulative function, and its strict tail are
\begin{align}
 W&=A_0+\sum_{j=1}^3p_j(c_j-c_{j-1})
    =\frac{2346784337}{1250000000},\label{eq:mass}\\
 F(h)&=A_0+\sum_{j=1}^3p_j[h-c_{j-1}]_0^{c_j-c_{j-1}}
       \quad(h\ge0),\label{eq:cdf}\\
 D(h)&=W-F(h)=\mu((h,+\infty))\quad(h\ge0).\label{eq:tail}
\end{align}
In particular, $W<2$, $A_0+p_{\max}<2$, $A_0<W/2$, and $p_{\max}<2$.
Set
\begin{equation}\label{eq:horizons}
 q=\lceil ak\rceil,\qquad H=k+q,\qquad L_0=2k+q,\qquad \cL=L_0+1,
 \qquad r_i=k-i\quad(i\in[k]).
\end{equation}

\subsection{The fractional cover}

For $(v,i)\in V\times[k]$, call $A_{v,i}=B(v,r_i)$ the original ball
with center $v$ and type $i$, and introduce a variable $x_{v,i}\ge0$.
For a ball $P\subseteq V$, put
\[
 T_x(P)=\sum_{v,i:A_{v,i}\cap P\ne\varnothing}x_{v,i},\qquad
 C_x(P)=\sum_{v,i:P\subseteq A_{v,i}}x_{v,i}.
\]
The linear feasibility problem is
\begin{align}
 \sum_{v\in V}x_{v,i}&=1 &&(i\in[k]),\label{eq:lp-type}\\
 \sum_{v,i:w\in A_{v,i}}x_{v,i}&\ge1 &&(w\in V),\label{eq:lp-cover}\\
 T_x(B(w,t))+C_x(B(w,t))&\ge2
       &&(w\in V, t\in\Z_{\ge0}, t\le L_0).
       \label{eq:lp-probe}
\end{align}
The balls in~\eqref{eq:lp-probe} are called probes.
A feasible solution of this system is called a fractional cover.

\begin{lemma}\label{lem:lp-validity}
If $b(G)\le k$, then~\eqref{eq:lp-type}--\eqref{eq:lp-probe} is
feasible. The feasibility problem has $nk$ variables, $O(nk)$
constraints, and a polynomial-time exact rational solution procedure.
\end{lemma}
\begin{proof}
Pad a covering sequence of length at most $k$ to length $k$, increasing
the earlier radii accordingly, and use its zero-one values for $x$.
For each probe $P$, either some chosen ball contains $P$, contributing
two to $T_x(P)+C_x(P)$, or at least two chosen balls meet $P$. These are
the only possibilities because the chosen balls cover $P$.
The constraints follow.
All ball intersections and inclusions are computable from graph distances.
Every coefficient is zero, one, or two. Exact rational feasibility of a
linear system with rational coefficients is solvable in polynomial time
in its encoding length~\cite{Khachiyan}; a feasible rational point of
polynomial encoding length can be returned.
\end{proof}

\subsection{Proposal labels and trials}

For $h\in\{0,\ldots,H\}$ define
\begin{equation}\label{eq:discrete-masses}
 \omega_0=A_0,\qquad
 \omega_h=\mu(((h-1)/k,h/k])\quad(1\le h\le H).
\end{equation}
Since $H/k\ge c_3$, these masses sum to $W$.
For a feasible $x$, a proposal label is a triple $p=(v,i,h)$ with
$\omega_hx_{v,i}>0$. Write $\mathcal P$ for the finite set of these labels.
Define its mass, raw ball, proposal ball, and raw
deadline by
\begin{equation}\label{eq:proposal-label}
 \lambda_p=\omega_hx_{v,i},\qquad
 A_p^0=B(v,r_i+h),\qquad A_p=B(v,r_i+h+1),\qquad
 t_p=H+i-h.
\end{equation}
Thus $t_p\in[L_0]$, the raw radius is $L_0-t_p$, and the proposal
radius is $\cL-t_p\le L_0$.

Put $N_0=10^{12}$ and $\delta=3/N_0$. For every $t\in[L_0]$, create
$N_0$ trials, each choosing label $p$ with probability $\lambda_p/N_0$
for each $p$ with $t_p=t$, and otherwise choosing a null outcome denoted
by $\bot$. All trials are independent. A label chosen by a trial is a
selected proposal. The intensity of a family of labels is the sum of
their masses; it is also the sum, over all trials, of their selection
probabilities.

The probabilities are well-defined. Indeed, with
$\rho_t=\sum_{p:t_p=t}\lambda_p$, at most one value of $i$ contributes
$\omega_0$, and every other $\omega_h$ is at most $p_{\max}/k$.
Using~\eqref{eq:lp-type} gives
\begin{equation}\label{eq:trial-bounds}
 \rho_t\le A_0+p_{\max}<2,\qquad
 \Prob(\text{a given trial is nonnull})\le\delta<\tfrac12,
 \qquad \sum_p\lambda_p=Wk<2k.
\end{equation}

A proposal is proper if $A_p$ is a proper subset of its connected
component. A proposal is a component proposal if $A_p$ is the whole
component $C$. For every component $C$, fix a center
$\kappa(C)\in\argmin_{v\in C}\ecc_C(v)$, resolving equal values by the
vertex ordering, and put $c(C)=\rad(C)$.
A component proposal is represented by $B(\kappa(C),c(C))$ and given
deadline $\cL-c(C)$. A proper proposal keeps its raw deadline.
Write $d_p$ for the resulting deadline and call $\cL-d_p$ its
advertised radius. This representation of a
component proposal preserves its set and satisfies $d_p\ge t_p$.

Two proposal labels belong to the same equality class if their proposal
balls are equal as vertex sets. A selected proposal set is
inclusion-maximal if no selected proposal set strictly contains it.
Retain one proposal from each inclusion-maximal selected proposal set. Among
selected proposals with the same set, prefer the largest $d_p$, then use
a fixed total order of their trial and label indices. This makes the
retained proposals uniquely determined. For $j\in[\cL]$, let $T_j$ be
the number of retained proposals whose deadline is at most $j$.
We will also use distributions obtained by replacing specified label
outcomes by $\bot$; their remaining trials stay independent.

\subsection{An intensity bound from the probes}

For $r\ge0$, define
\begin{equation}\label{eq:g-definition}
 g(r)=\min\left\{W,
  2W-\max_{0\le t\le r}\bigl(F(r-t)+F(r+t)\bigr)\right\}.
\end{equation}
This function is nonnegative, nonincreasing, continuous, and piecewise
affine. It has Lipschitz constant at most $2p_{\max}$, meaning that
$|g(r)-g(r')|\le2p_{\max}|r-r'|$. Also $g(r)=0$ for $r\ge c_3$ and
$g(0)=W$. To verify these properties and to evaluate $g$ exactly, use
\begin{equation}\label{eq:g-finite}
 \max_{0\le t\le r}\bigl(F(r-t)+F(r+t)\bigr)
 =\max\left(\{2F(r)\}\cup
 \{F(z)+F(2r-z):z\in\{c_0,c_1,c_2,c_3\},\ z\le2r\}\right).
\end{equation}
For completeness, the sum $F(y)+F(2r-y)$ is affine between its finitely
many breakpoints on $[0,2r]$. Its endpoint and breakpoint values are the
values on the right, using symmetry about $r$.
Each active value is nondecreasing with slope at most $2p_{\max}$.
A newly admissible breakpoint has the endpoint value when it appears.
Taking the maximum and then the minimum in~\eqref{eq:g-definition}
proves the asserted continuity, monotonicity, and Lipschitz bound.

\begin{lemma}\label{lem:absorption}
For every $v\in V$ and integer $0\le R\le L_0$, the raw proposals
containing $B(v,R)$ have intensity at least $g(R/k)$.
\end{lemma}
\begin{proof}
Call an original ball universal for this target if it contains $B(v,R)$,
and let $U_0$ be the sum of its $x$ values. Each unit of universal mass
contributes all $W$ enlargement mass to raw containment. If $U_0\ge1$,
the conclusion follows from $g\le W$. If $U_0<1$, remove universal
balls and divide remaining masses by $1-U_0$. Coverage of $v$ remains
at least one. For every probe $B(v,t)$ with $t\le R$, each removed ball
contributed twice its mass to~\eqref{eq:lp-probe}, so these probe
inequalities remain valid after normalization. We first prove the
conclusion for this normalized system without universal balls.

Let $C\ge1$ be its mass of original balls containing $v$, called near
balls. All other original balls are called far balls. For a near ball,
let $\tau$ be the minimum distance from $v$ to a vertex outside the ball.
Since the ball is not universal, $1\le\tau\le R$ and it contains
$B(v,\tau-1)$. For any vertex of $B(v,R)$, follow a shortest path from
$v$ for its first $\tau-1$ edges, or stop earlier at the vertex itself.
It follows that enlargement $R-\tau+1$ suffices to contain $B(v,R)$.
A far ball meeting $B(v,\tau)$ contains the target after enlargement
$R+\tau$, by the triangle inequality.

For integer $t\in[0,R]$, let $Q(t)$ be near mass with exit distance
greater than $t$, and let $N(t)$ be far mass meeting $B(v,t)$. The probe
inequality reads
\begin{equation}\label{eq:near-far}
 C+Q(t)+N(t)\ge2.
\end{equation}
Every near ball contains the target after any enlargement strictly
greater than $R$. From~\eqref{eq:g-definition} with $t=0$,
\begin{equation}\label{eq:near-tail}
 W-F(R/k)\ge\tfrac12g(R/k).
\end{equation}
If $C\ge2$, two units of near mass suffice by~\eqref{eq:near-tail}.

Suppose $1\le C<2$. Select near mass $2-C$ with the largest exit
distances, allowing a partial mass at a boundary value. Its mass with
exit distance at most $t$ is
$\max\{0,2-C-Q(t)\}\le N(t)$ by~\eqref{eq:near-far}.
Match this selected near mass to far mass, requiring that a near ball
with exit distance $\tau$ is matched only to a far ball meeting
$B(v,\tau)$. Such a fractional matching exists: process exit distances
in increasing order and use previously unused eligible far mass.
Eligibility sets are nested, and the preceding cumulative inequality
ensures enough eligible mass at every step.

For each matched unit, the two raw containment contributions are at least
\begin{equation}\label{eq:paired-tails}
 \mu(((R-\tau)/k,+\infty))+
 \mu(((R+\tau)/k,+\infty))\ge g(R/k).
\end{equation}
Strictness of the tails and upward rounding in
\eqref{eq:discrete-masses} ensure the integer enlargements stated above.
The matched mass contributes $(2-C)g(R/k)$.
Remaining near mass is $2(C-1)$ and contributes at least
$(C-1)g(R/k)$ by~\eqref{eq:near-tail}. These contributions total $g(R/k)$.
Restoring the universal mass gives at least
$U_0W+(1-U_0)g(R/k)\ge g(R/k)$.
When $R=0$, every original ball covering $v$ is universal, so the first
case applies directly.
\end{proof}

\begin{lemma}\label{lem:strict-growth}
If a target ball $A=B(v,R)$ is proper in its component and $R\le L_0$,
the proposal balls strictly containing $A$ have intensity at least
$g(R/k)$.
\end{lemma}
\begin{proof}
For $S\subseteq V$, define its closed neighborhood by
$N[S]=\{w\in V:d_G(w,S)\le1\}$, where
$d_G(w,S)=\min_{z\in S}d_G(w,z)$ and $\min\varnothing=+\infty$.
Graph balls satisfy $N[B(u,r)]=B(u,r+1)$ for integer $r\ge0$.
Thus every raw ball containing $A$ gives a proposal ball containing
$N[A]$. A nonempty proper subset of a connected component has an edge
to its complement, so $A\subsetneq N[A]$.
Apply Lemma~\ref{lem:absorption}.
\end{proof}

\subsection{Bounds for expected counts}

Define for $z\ge0$ and $r\ge0$
\begin{equation}\label{eq:phi-q}
 \phi(z)=\int_0^1e^{-tz}\,dt
 =\begin{cases}(1-e^{-z})/z,&z>0,\\1,&z=0,\end{cases}
 \qquad q_P(r)=e^{-g(r)}.
\end{equation}
The subscript $P$ in $q_P$ is a fixed label for proper proposals.
The function $\phi$ is nonincreasing, $e^{-z}\le\phi(z)\le1$, and
$|\phi'(z)|\le1/2$. In particular, $q_P$ and $\phi\circ g$ are
nondecreasing, their Lipschitz constants are at most $2p_{\max}$ and
$p_{\max}$, and $q_P(r)\le\phi(g(r))$.

\begin{lemma}\label{lem:count-bounds}
Before any deletion of label outcomes, the expected number of retained
proper proposals of deadline at most $j$ is at most
\begin{equation}\label{eq:proper-count}
 e^\delta\sum_{\substack{p\text{ proper}\\d_p\le j}}
         \lambda_pq_P((\cL-t_p)/k).
\end{equation}
For a component $C$, let $\lambda_C$ be the intensity of component
proposals with set $C$. The probability that its component proposal is
retained is at most
\begin{equation}\label{eq:component-count}
 e^\delta\lambda_C\phi(g(c(C)/k)).
\end{equation}
The same conclusions hold after deleting labels if no containing label
of a remaining target has been deleted; components entirely deleted
can be omitted.
\end{lemma}
\begin{proof}
Consider an equality class of proper proposal sets, with strict-superset
intensity $\lambda_s$. If this class is retained by $j$, some outcome
in the class with deadline at most $j$ is selected and every strict
superset is avoided. Conditional on selecting that outcome in its trial,
independence of the other trials bounds the avoidance probability by
$e^{-\lambda_s+\delta}$. The omitted trial has strict-superset
probability at most $\delta$. For each representation of the class,
Lemma~\ref{lem:strict-growth} gives
$\lambda_s\ge g((\cL-t_p)/k)$. A union bound over the eligible selected
outcomes, followed by summing over equality classes, proves
\eqref{eq:proper-count}.

For a component, write $s_\xi$ for its selection probability in trial
$\xi$, so $\sum_\xi s_\xi=\lambda_C$ and $s_\xi\le\delta$.
Its probability of at least one selection is
\[
 1-\prod_\xi(1-s_\xi)
 =\int_0^1\sum_\xi s_\xi\prod_{\zeta\ne\xi}(1-t s_\zeta)\,dt
 \le e^\delta\lambda_C\phi(\lambda_C).
\]
If $c(C)\le L_0$, Lemma~\ref{lem:absorption} at $\kappa(C)$ gives
$\lambda_C\ge g(c(C)/k)$. If $c(C)>L_0$, no component proposal exists
and $g(c(C)/k)=0$. Monotonicity of $\phi$ proves
\eqref{eq:component-count}. The deletion assertion follows because the
same required intensities remain available.
\end{proof}



\section{Analytic bounds and the numerical inequality}\label{sec:analytic}

For a radius parameter $r\ge0$, let $\nu_r$ be the translate of $\mu$ by $r$:
$\int f(x)\,d\nu_r(x)=\int f(r+h)\,d\mu(h)$.
Define the finite measure $\Lambda$ by
\begin{equation}\label{eq:aggregate-measure}
 \int f(x)\,d\Lambda(x)=\int_0^1\int f(x)\,d\nu_r(x)\,dr.
\end{equation}
It has total mass $W$, assigns zero mass outside $[0,1+c_3]$, and has a density
\begin{equation}\label{eq:aggregate-density}
 f_\Lambda(x)=A_0\ind_{x\in[0,1]}
 +\sum_{j=1}^3p_j
   \pos{\min\{1,x-c_{j-1}\}-\max\{0,x-c_j\}}.
\end{equation}
A density is a nonnegative function whose integral over each interval
equals its $\Lambda$ measure. Its values at the finitely many endpoints
do not affect the integral. Integrating the translated atom and the
translated density intervals in~\eqref{eq:aggregate-measure}
proves~\eqref{eq:aggregate-density}.
In particular, $f_\Lambda\le A_0+p_{\max}<2$.

For $e\ge\theta\ge0$, define
\begin{equation}\label{eq:reward}
 \Phi_{\theta,e}(x)=
 \begin{cases}
  0,&x<\theta,\\
  q_P(x),&\theta\le x<e,\\
  \phi(g(e)),&x\ge e.
 \end{cases}
\end{equation}
Call $e$ the cut value of this function. Define
\begin{align}
 K(r,\theta)&=\sup_{e\ge\theta}\int\Phi_{\theta,e}(x)\,d\nu_r(x),
       &M_N(\theta)&=\int_0^1K(r,\theta)\,dr,\label{eq:new-envelope}\\
 M_O(\theta)&=\int_{[\theta,+\infty)}e^{-D(x)}\,d\Lambda(x),
       &M(\theta)&=\min\{M_N(\theta),M_O(\theta)\}.
       \label{eq:old-envelope}
\end{align}
The subscripts are fixed names for the two bounds, not parameters.

\begin{lemma}\label{lem:kernel-properties}
For fixed $\theta$, $K(r,\theta)$ is nondecreasing in $r$ and lies in
$[0,W]$. For $0\le\theta'\le\theta$,
\begin{equation}\label{eq:kernel-threshold}
 0\le K(r,\theta')-K(r,\theta)
       \le\nu_r([\theta',\theta)).
\end{equation}
Each of $M_N,M_O,M$ is nonincreasing, takes values in $[0,W]$, and
satisfies
\begin{equation}\label{eq:envelope-increment}
 0\le M_*(\theta')-M_*(\theta)
       \le\Lambda([\theta',\theta))\le2(\theta-\theta'),
 \qquad M_*\in\{M_N,M_O,M\}.
\end{equation}
\end{lemma}
\begin{proof}
The function in~\eqref{eq:reward} is nondecreasing in $x$ and lies in
$[0,1]$: for $x<e$, $q_P(x)\le\phi(g(x))\le\phi(g(e))$.
Translation of $\nu_r$ proves monotonicity in $r$ and the bound $W$.
For~\eqref{eq:kernel-threshold}, a cut $e\ge\theta$ loses at most the
mass of $[\theta',\theta)$ when the threshold increases. If instead
$\theta'\le e<\theta$, replace $e$ by $\theta$. This can only increase
the reward on $[\theta,+\infty)$, and again the lost reward is at most
the indicated interval mass. Taking suprema proves the assertion.
Integrating in $r$ proves it for $M_N$. For $M_O$ it follows because
$0\le e^{-D}\le1$. The minimum of two nonincreasing functions with the
same increment bound has that increment bound. The density bound then
proves the last inequality.
\end{proof}

For $s\in[0,2]$, put
\begin{equation}\label{eq:analytic-capacity}
 A(s)=\min_{s\le t\le2}\{a+t-M(2-t)\}.
\end{equation}
The minimum exists by~\eqref{eq:envelope-increment}, and $A$ is
nondecreasing. When $A>0$, define
\begin{equation}\label{eq:integral-J}
 J=\frac12\int_0^2\frac{e^{-g(1-s/2)}}{A(s)}\,ds.
\end{equation}
Its integrand is nonincreasing: the numerator is nonincreasing in $s$
and the denominator is positive and nondecreasing.

\begin{lemma}[Exact numerical inequality]\label{lem:numerical}
For the constants in~\eqref{eq:profile},
\begin{align}
 A(s)&\ge\frac{54337893542773910}{524288000000000000}
           >\frac1{10} &&(0\le s\le2),\label{eq:cert-capacity}\\
 J&\le\frac{65533759887597263}{65536000000000000}
       =0.9999658185973703460693359375
       <\frac{999999}{1000000}.\label{eq:cert-J}
\end{align}
\end{lemma}
\begin{proof}
Appendix~\ref{app:certificate} gives finite rational upper and lower
bounds for the defining integrals, the supremum in $K$, and the minimum
in $A$. It proves the direction of every bound and specifies the integer
evaluation of those bounds. With $65536$ intervals for the radius parameter, $65536$ cut
intervals, and $524288$ primitive intervals per unit, the evaluated
lower numerator for the uniform capacity is $54337893542773910$ with
denominator $524288000000000000$. The evaluated upper numerator for
$J$ is $131067519775194526$ with denominator $131072000000000000$.
Reduction of the latter fraction gives~\eqref{eq:cert-J}.
The complete programs performing this finite calculation, including
the rational profile, are part of this LaTeX source and are displayed
in that appendix.
\end{proof}

\section{Deterministic rounding of the proposals}\label{sec:rounding}

Throughout this section fix a rational $\varepsilon\in(0,1]$, and put
\begin{equation}\label{eq:tolerance}
 \sigma=\min\{10^{-7},\varepsilon/16\},\qquad
 K_0=\max\{10^{12},\lceil16/\varepsilon\rceil\}.
\end{equation}
We will specify which guesses use the construction of this section.

\subsection{Geometric grouping and guards}

Partition the integer radii $0,\ldots,\cL-1$ into a singleton group for
zero and the nonempty groups induced by
$[(1+\sigma)^h,(1+\sigma)^{h+1})$, for integers $h\ge0$.
An advertised radius $R$ corresponds to deadline $\cL-R$.
For each induced group of deadlines keep its largest deadline, even if
no positive-mass proposal has a deadline in the group. Let
$\mathcal J\subseteq[\cL]$ be this set of endpoints and
$\ell=|\mathcal J|$. In particular $\cL\in\mathcal J$.
For $j\in[\cL]$, let $j^+$ be the smallest endpoint at least $j$.
The width of every group gives
\begin{equation}\label{eq:group-width}
 0\le j^+-j\le w,\qquad w=\lceil\sigma\cL\rceil+2,
 \qquad \ell=O_\sigma(\log(n+1)).
\end{equation}
For the last bound, $\cL<3n$ whenever $k\ge K_0$; the finitely many
smaller guesses will use Theorem~\ref{thm:fpt}.

Put
\begin{equation}\label{eq:guard-eta}
 \eta=\frac{\sigma^2}{10^8\ell}.
\end{equation}
For a candidate proposal ball $A_p$, its containment downset is the
set of currently undeleted labels $q$ with $A_q\subseteq A_p$.
If any containment downset has intensity at least $\eta k$, record its
candidate ball as a guard and delete all labels in that downset.
Repeat until every candidate has downset intensity less than $\eta k$.
Write $\mathcal G$ for the list of guard balls and
$N_{\mathrm g}=|\mathcal G|$.

\begin{lemma}\label{lem:guards}
This procedure is polynomial-time and
\begin{equation}\label{eq:guard-count}
 N_{\mathrm g}<2/\eta.
\end{equation}
If a target set is not contained in a guard, none of its containing
proposal labels has been deleted. If a component proposal is deleted,
its whole component is contained in a guard.
\end{lemma}
\begin{proof}
Each iteration deletes at least $\eta k$ intensity, and disjoint deleted
families have total intensity at most $Wk<2k$. This proves the bound.
The number of iterations is also at most the number of labels, since
each iteration deletes a positive-mass label. Containments and sums
are polynomial-time computations. If a deleted proposal contains a
target, its guard contains that target. The component assertion is its
special case.
\end{proof}

Replace every deleted label outcome by $\bot$ in every trial. The
remaining trials are independent and obey~\eqref{eq:trial-bounds}.
For remaining candidates and for original balls not contained in guards,
Lemmas~\ref{lem:absorption}--\ref{lem:count-bounds} remain valid by
Lemma~\ref{lem:guards}.
All definitions of $T_j$ below refer to this modified distribution.

\subsection{Variance of a deadline count}

Expectations in this subsection refer to the independent trials after
guard deletion. Conditional expectations fix specified trial values
and average over the remaining trial values.

\begin{lemma}\label{lem:deletion-variance}
Let $X_1,\ldots,X_N$ be independent finite-valued random variables,
each with a designated null value. For a real function $Z$ of these
variables, let $Z^{(-t)}$ be obtained by replacing $X_t$ by its null
value. Then
\[
 \Var Z\le\sum_{t=1}^N\E(Z-Z^{(-t)})^2.
\]
\end{lemma}
\begin{proof}
Put $Y_t=\E[Z\mid X_1,\ldots,X_t]$, with $Y_0=\E Z$.
Successive differences are orthogonal: for $s<t$, conditioning on
$X_1,\ldots,X_{t-1}$ gives
$\E[(Y_s-Y_{s-1})(Y_t-Y_{t-1})]=0$.
Hence $\Var Z=\sum_t\E(Y_t-Y_{t-1})^2$.
Let $\E_tZ$ denote integration over $X_t$ with all other coordinates
fixed. Independence gives
$Y_t-Y_{t-1}=\E[Z-\E_tZ\mid X_1,\ldots,X_t]$.
For real numbers $y_i$ and nonnegative weights $p_i$ summing to one,
$(\sum_i p_iy_i)^2\le\sum_i p_iy_i^2$. Applying this inequality
to the conditional average bounds its squared expectation by
$\E(Z-\E_tZ)^2$.
For fixed other coordinates the mean minimizes squared error, so this
quantity is at most $\E(Z-Z^{(-t)})^2$. Sum over $t$.
\end{proof}

\begin{lemma}\label{lem:count-variance}
After deletion of the guard downsets, for every $j\in[\cL]$,
\begin{equation}\label{eq:count-variance}
 \Var T_j\le6k+4\eta k^2.
\end{equation}
\end{lemma}
\begin{proof}
Treat each possible nonnull outcome as a pair consisting of a trial and
a label. Order these outcomes by strict set inclusion; order equal sets
by the deadline and the fixed tie order. This is a strict partial order.
Its maximal selected outcomes are exactly the retained proposals.
An outcome above another in this order is called its dominator.

Let $N$ be the number of selected outcomes. Deleting a selected
nonmaximal outcome does not change any maximal outcome. If deleting a
maximal outcome $e$ creates $a_e$ new maximal outcomes, each created
outcome had $e$ as its unique selected dominator. Set $a_e=0$ for every
possible outcome that is not selected and maximal; sums over $e$ below
range over all possible nonnull outcomes. These created-outcome
sets are disjoint for different $e$, and consequently
\begin{equation}\label{eq:exposure-disjoint}
 \sum_e a_e\le N.
\end{equation}
The squared change in a prefix count is at most $1+a_e^2$.

Call $(e,f)$ an exposed pair if $f$ becomes maximal on deleting $e$.
Fix such a pair of positive probability.
The event defining this pair specifies that $e$ and $f$ are selected in
distinct trials and that every other trial avoids the dominators of
$f$. Transitivity also makes $e$ maximal in this event.
This conditioning is separate for each other trial. Its avoidance
probability is at least $1-\delta$, so the probability of any allowed
nonnull outcome increases by at most $(1-\delta)^{-1}\le2$.
Every outcome exposed by deleting $e$ is in the containment downset of
$e$. Its unconditioned intensity is less than $\eta k$.
The two fixed trials and the remaining trials therefore give
\[
 \E[a_e\mid(e,f)\text{ is an exposed pair}]\le2+2\eta k.
\]
Writing $a_e^2$ as the sum of $a_e$ over its exposed outcomes yields
\[
 \E\sum_ea_e^2\le(2+2\eta k)\E\sum_ea_e
                    \le(2+2\eta k)\E N.
\]
Lemma~\ref{lem:deletion-variance} now gives
\[
 \Var T_j\le\E N+\E\sum_e a_e^2.
\]
Since
$\E N\le Wk<2k$, the result is
$(3+2\eta k)\E N\le6k+4\eta k^2$.
\end{proof}

\subsection{Discrete expected counts}

Put $\mu_j=\E T_j$. These expectations depend on the modified trial
distribution after guards. For $j\in[\cL]$ put
$\theta_j=\max\{0,(L_0-j)/k\}$.

\begin{lemma}\label{lem:discrete-means}
For every $q\le j\le\cL$,
\begin{equation}\label{eq:discrete-mean-bound}
 \mu_j\le kM(\theta_j)+100+100k/N_0.
\end{equation}
For $k\ge10^{12}$ and $1\le j<q$,
\begin{equation}\label{eq:early-mean}
 \mu_j\le\frac{p_{\max}j(j+1)}{2k}\le j.
\end{equation}
\end{lemma}
\begin{proof}
First use the threshold $\theta=(\cL-j)/k\ge0$.
For a center $v$ in a component $C$, put
$e_v=\ecc_C(v)/k$. For a label $p=(v,i,h)$, put
$x=(r_i+h+1)/k$, its normalized proposal radius. Its proposal is proper
exactly when $x<e_v$.
If this proper proposal is counted by $j$, then $x\ge\theta$ and
the coefficient of its mass in~\eqref{eq:proper-count} is
$e^\delta q_P(x)$.
If its component proposal is counted by $j$, then
$c(C)/k\ge\theta$ and $e_v\ge\theta$. Since $c(C)/k\le e_v$,
the coefficient of its mass in~\eqref{eq:component-count} is at most
$e^\delta\phi(g(e_v))$. If $e_v<\theta$, proposals centered at $v$
contribute no count by $j$. Otherwise, the function
$e^\delta\Phi_{\theta,e_v}$ bounds the coefficients in the two
expected-count bounds, including a zero coefficient when a proposal
is not counted by $j$.

For $h>0$ in the continuous enlargement measure, upward discretization
replaces $h$ by $\lceil kh\rceil/k$; for the atom at zero it leaves
$h=0$. In both cases, the normalized proposal radius is at most
$r_i/k+h+2/k$. Monotonicity of $\Phi$ in its radius argument,
followed by the definition of $K$ and~\eqref{eq:lp-type}, gives
\[
 \mu_j\le e^\delta\sum_{i=1}^k K((k-i+2)/k,\theta).
\]
For any nondecreasing function $f:[0,1]\to[0,W]$, its right Riemann
sum satisfies
\[
 \sum_{i=1}^kf(i/k)\le k\int_0^1f(r)\,dr+W.
\]
Apply this inequality with $f(r)=K(r,\theta)$. Replacing the first
summand $K(1/k,\theta)$ by $K(1+1/k,\theta)\le W$ adds at most $W$.
It follows that
\begin{equation}\label{eq:new-discrete-bound}
 \mu_j\le e^\delta\{kM_N(\theta)+2W\}
           \le kM_N(\theta_j)+100+100k/N_0.
\end{equation}
Here $\theta\ge\theta_j$, $W<2$, $M_N\le W$, and
$e^\delta\le1+6/N_0$ have been used. This proof also applies when
$j=\cL$ and $e_v=0$.

For a second bound, map each retained set to a selected outcome with
that set and minimum raw deadline. Its raw deadline is at most its
retained deadline, so a set counted
by $j$ is mapped to a raw radius at least $k\theta_j$.
For a raw ball $B(v,S)$, every original ball covering $v$ and enlarged
by more than $S$ gives a containing raw ball with a strictly earlier raw
deadline. Its proposal ball also contains $B(v,S+1)$.
Such earlier outcomes have intensity at least $D(S/k)$ by
\eqref{eq:lp-cover}. A selected earlier containing set would either
contradict maximality or contradict the minimum raw deadline.
All these outcomes are in earlier trials, so independence bounds the
probability of such an outcome being selected and eligible for this map
by its selection probability times $e^{-D(S/k)}$. Different retained
sets map to different outcomes, so summing bounds $\mu_j$. Guard deletion preserves the required earlier outcomes
by Lemma~\ref{lem:guards}.

Define a finite measure $\Lambda_k$ by
\[
 \int f(x)\,d\Lambda_k(x)
 =\frac1k\sum_{i=1}^k\sum_{h=0}^H
       \omega_h f((r_i+h)/k).
\]
This is the total measure of normalized raw radii, divided by $k$;
the equality follows from~\eqref{eq:lp-type}.
For each $i$, compare $r_i/k+\lceil kh\rceil/k$ with $r+h$ for
$r\in[r_i/k,(r_i+1)/k]$ and $h\ge0$.
The difference in absolute value is at most $1/k$, including $h=0$.
Integrate this comparison in $r$ and with respect to $\mu(h)$.
Since $\Lambda$ has density less than two, it gives
\[
 \Lambda_k([x,+\infty))\le\Lambda([x,+\infty))+2/k
 \quad\text{for every }x\ge0.
\]
The function $x\mapsto\ind_{x\ge\theta_j}e^{-D(x)}$ is nondecreasing,
with values in $[0,1]$. Integrating its increments against the preceding
tail inequality gives
\begin{equation}\label{eq:old-discrete-bound}
 \mu_j\le kM_O(\theta_j)+2.
\end{equation}
This increment integration follows first for nondecreasing step
functions by summing their jumps, and then for the displayed function
by monotone step approximations.
Taking the better of~\eqref{eq:new-discrete-bound} and
\eqref{eq:old-discrete-bound} proves~\eqref{eq:discrete-mean-bound}.

If $j<q$, a nonzero label at deadline $t\le j$ has $i\le t$ and its
positive enlargement mass is at most $p_{\max}/k$. Thus
$\rho_t\le p_{\max}t/k$. Retention and the representation of component
proposals only delete outcomes or delay their deadlines, giving the first inequality
in~\eqref{eq:early-mean}. The last one follows from $j<q$,
$a<0.12$, $p_{\max}<2$, and $k\ge10^{12}$.
\end{proof}


\subsection{Capacities and a joint potential}

Assume from now on that
\begin{equation}\label{eq:large-regime}
 k\ge K_0,\qquad k\ge\frac{10^{11}\ell}{\sigma^3}.
\end{equation}
For $q\le j\le L_0$, equations~\eqref{eq:analytic-capacity},
\eqref{eq:cert-capacity}, and~\eqref{eq:discrete-mean-bound} give
\[
 j-\mu_j\ge(q-ak)+kA(0)-100-100k/N_0>0.
\]
The same calculation at $j=\cL$ has one additional position.
Together with~\eqref{eq:early-mean}, this proves
\begin{equation}\label{eq:means-feasible}
 \mu_j\le j\qquad(j\in[\cL]).
\end{equation}
Define integers
\begin{equation}\label{eq:U-B}
 U_j=\lceil\mu_{j^+}+\sigma k\rceil,\qquad
 B=w+\lceil\sigma k\rceil+2,\qquad D_i=q+2i.
\end{equation}
The deadline $D_i$ permits a radius-$2r_i$ repair in the base horizon,
since $\cL-D_i=2r_i+1$. Put
\begin{equation}\label{eq:integer-capacities}
 a_i=\min_{D_i\le j\le\cL}(B+j-U_j)\qquad(i\in[k]).
\end{equation}
These are nondecreasing integers. The width bound gives
$B+j-U_j\ge j^+-\mu_{j^+}+1$.
For $j\ge D_i$, its endpoint satisfies
$\theta_{j^+}\le2-2i/k$. Using~\eqref{eq:analytic-capacity} and
\eqref{eq:discrete-mean-bound}, including the final endpoint with
$\theta_{\cL}=0$, proves
\begin{equation}\label{eq:capacity-lower}
 a_i\ge kA(2i/k)-100-100k/N_0>0.
\end{equation}
Also~\eqref{eq:means-feasible} and~\eqref{eq:group-width} imply
\begin{equation}\label{eq:U-feasible}
 U_j\le B+j.
\end{equation}

For each original ball $A_{v,i}$, let $I_{v,i}$ be the indicator that it
is not contained in a guard and is not contained in any selected
proposal ball. If it is not contained in a guard, independence and
Lemma~\ref{lem:absorption} bound its noncontainment probability by
$e^{-g(r_i/k)}$. Indeed, none of its containing proposals was deleted
by the guards. Linearity of expectation and~\eqref{eq:lp-type} give
\begin{equation}\label{eq:residual-expectation}
 \E\sum_{v,i}\frac{x_{v,i}I_{v,i}}{a_i}
 \le\sum_{i=1}^k\frac{e^{-g(1-i/k)}}{a_i}
 \le(1+2000/k+2000/N_0)J.
\end{equation}
Here is the last estimate in detail. From~\eqref{eq:capacity-lower}
and $A>1/10$, the relative capacity error is at most
$d=1000/k+1000/N_0<1/2$.
The inequality $(1-d)^{-1}\le1+2d$ bounds the reciprocals.
The right Riemann sum for the nonincreasing integrand in
\eqref{eq:integral-J} is at most its integral, giving
\eqref{eq:residual-expectation}.

Consider the following nonnegative function of the trial outcomes:
\begin{equation}\label{eq:potential}
 \Psi=\sum_{v,i}\frac{x_{v,i}I_{v,i}}{a_i}
        +\sum_{j\in\mathcal J}
            \frac{(T_j-\mu_j)^2}{\sigma^2k^2}.
\end{equation}
All its coefficients are rational and fixed before the outcomes are
chosen. Lemma~\ref{lem:count-variance} and
\eqref{eq:residual-expectation} imply
\begin{align}
 \E\Psi
 &\le(1+2000/k+2000/N_0)J
            +\frac{6\ell}{\sigma^2 k}+\frac{4\ell\eta}{\sigma^2}
            \label{eq:potential-expectation}\\
 &<(1+4\cdot10^{-9})\frac{999999}{1000000}
                      +6\cdot10^{-18}+4\cdot10^{-8}<1.
 \nonumber
\end{align}
The second line uses~\eqref{eq:large-regime},
\eqref{eq:tolerance},~\eqref{eq:guard-eta}, and
Lemma~\ref{lem:numerical}.

\begin{lemma}\label{lem:conditional-computation}
The unconditional means $\mu_j$ and every conditional expectation of
$\Psi$ obtained by fixing trial outcomes are exactly computable in
polynomial time, with rational numbers of polynomial encoding length.
Consequently a deterministic polynomial-time choice of outcomes
satisfies $\Psi<1$.
\end{lemma}
\begin{proof}
For an original ball not contained in a guard, noncontainment is the
event that every trial avoids its containing labels. Its probability is
a product of rational avoidance probabilities.

For completeness, consider the partial order in the proof of
Lemma~\ref{lem:count-variance}. Let $o=(\xi,p)$ be a possible nonnull
outcome in trial $\xi$, and let $\pi_{\xi,p}$ be its current probability.
Write $o'\succ o$ when $o'$ dominates $o$.
The probability that $o$ is selected and maximal is
\[
 \pi_{\xi,p}\prod_{\zeta\ne\xi}
       \left(1-\sum_{q:(\zeta,q)\succ o}\pi_{\zeta,q}\right).
\]
For two outcomes in different trials that are incomparable, their joint
maximality probability is the product of their selection probabilities
and the avoidance probabilities of the union of their dominators in
all other trials. Comparable outcomes cannot both be maximal.
Distinct outcomes of one trial cannot both be selected.
These formulas compute first and second moments of each $T_j$.
Fixed trials have point-mass distributions, so the same formulas compute
conditional moments. In particular,
$\E[(T_j-\mu_j)^2\mid\text{fixed outcomes}]$ follows from its conditional
first and second moments and the fixed unconditional $\mu_j$.

The number of trials is $N_0L_0$, and there are polynomially many
possible labels per trial. Thus the numbers of products and pairs are
polynomial. Each product has polynomially many rational factors of
polynomial encoding length. Their products, sums, and squares also have
polynomial encoding length. This proves the computation assertion.

Fix the trials successively. At every step the current conditional
expectation is the probability-weighted average of the conditional
expectations after the possible next choices. Some positive-probability
choice therefore does not increase it. Select the first such choice.
At termination $\Psi\le\E\Psi<1$.
\end{proof}

For these deterministically selected outcomes, every term of
\eqref{eq:potential} is less than one. In particular,
\begin{equation}\label{eq:deterministic-guarantees}
 \sum_{v,i}\frac{x_{v,i}I_{v,i}}{a_i}<1,\qquad
 |T_j-\mu_j|<\sigma k\quad(j\in\mathcal J).
\end{equation}
Monotonicity of $T_j$ gives $T_j\le U_j$ for every integer
$j\in[\cL]$.

\subsection{Residual coverage and assignment of rounds}

Let $U$ be the vertices outside the union of the guards and the retained
proposal balls, and put $z_{v,i}=x_{v,i}I_{v,i}$. The retained proposal
sets have the same union as all selected proposal sets: every selected
set is contained in an inclusion-maximal selected set.
If an original ball contains a vertex of $U$, it is not contained in a
guard or a selected proposal. Thus
\[
 \sum_{v,i:w\in A_{v,i}}z_{v,i}\ge1\qquad(w\in U).
\]
The first inequality of~\eqref{eq:deterministic-guarantees} and
Lemma~\ref{lem:cumulative-rounding}, with radii $r_i$ and capacities
$a_i$, give a cover of $U$ by balls of radius $2r_i$, which we call
repair balls of type $i$. Their numbers satisfy all the cumulative
count bounds in that lemma.

Reserve $N_{\mathrm g}$ initial rounds for the guards. In the following
timeline, of length $B+\cL$, give each retained proposal deadline
$B+d_p$. There are $T_j\le U_j\le B+j$ such jobs due by $B+j$, and
there are no jobs due in $[B]$. Lemma~\ref{lem:deadlines} therefore
places these jobs as late as possible. Its formula gives at least
\[
 \min_{D_i\le j\le\cL}(B+j-T_j)\ge a_i
\]
free positions by $B+D_i$.
Place the repair balls in increasing type order, at the earliest free
positions. At most $a_i$ repairs have type at most $i$, so every type-$i$
repair is placed by $B+D_i$.

The resulting horizon is
\begin{equation}\label{eq:final-horizon}
 L=N_{\mathrm g}+B+\cL.
\end{equation}
A retained proposal placed by $N_{\mathrm g}+B+d_p$ has propagation
radius at least $\cL-d_p$, exactly its advertised radius.
A type-$i$ repair placed by $N_{\mathrm g}+B+D_i$ has radius at least
$2r_i+1$, which contains its radius-$2r_i$ repair ball.
Every guard has radius at most $L_0<\cL$ and is placed in the first
$N_{\mathrm g}$ rounds, leaving at least $B+\cL$ propagation rounds.
The guards, retained proposals, and repairs cover $V$.
Lemma~\ref{lem:legalization} gives a legal sequence of length at most $L$.

\subsection{Running time and approximation factor}

All operations for a guess satisfying~\eqref{eq:large-regime} are
polynomial-time bit operations. Graph distances, ball sets, and
containments have polynomial size. The fractional cover is obtained by
Lemma~\ref{lem:lp-validity}. The number of label triples is at most
$nk(H+1)$, and $N_0$ is an absolute constant.
Guard deletion, grouping, the computations of
Lemma~\ref{lem:conditional-computation}, and the rounding of
Lemma~\ref{lem:cumulative-rounding} are polynomial-time operations.
Geometric endpoints can be compared with integers by exact powers of
the rational $1+\sigma$; their number and bit lengths are polynomial
for fixed $\sigma$. No evaluations of real exponential functions are
needed by the input-dependent algorithm: $M$, $A$, $g$, and $J$ occur
in its analysis, while $\mu_j$, $U_j$, $a_i$, and all conditional
expectations are computed as rational numbers.

If~\eqref{eq:large-regime} fails, run Theorem~\ref{thm:fpt} instead.
If $k<K_0$, its exponential factor is at most $4^{K_0}$, a constant
depending only on $\varepsilon$. Otherwise
$k<10^{11}\ell/\sigma^3=O_\varepsilon(\log(n+1))$, and
$4^k=(n+1)^{O_\varepsilon(1)}$. Thus the fallback is polynomial-time
for every guess where it is used.

For a large-regime success, equations~\eqref{eq:U-B},
\eqref{eq:group-width}, and~\eqref{eq:final-horizon} give
\begin{align}
 L&\le(2+a+(3+a)\sigma)k+8+2\sigma+N_{\mathrm g},
       \label{eq:length-overhead}\\
 \frac{N_{\mathrm g}}k
 &<\frac{2\cdot10^8\ell/\sigma^2}
          {10^{11}\ell/\sigma^3}
   =0.002\sigma. \label{eq:guard-ratio}
\end{align}
The coefficient $3+a+0.002$ is less than four. Therefore
\[
 \frac Lk<2+a+4\sigma+\frac9k
       \le2+a+\frac{\varepsilon}{4}+\frac{9\varepsilon}{16}
       <2+a+\varepsilon.
\]
A fallback success has length at most $2k-1$ and obeys the same
approximation bound.

\begin{proof}[Proof of Theorem~\ref{thm:poly}]
For a fixed rational $\varepsilon\in(0,1]$, test $k=1,\ldots,n$.
Use the fallback when~\eqref{eq:large-regime} fails; otherwise solve
the fractional-cover problem and reject an infeasible guess. At a
feasible large-regime guess, run the construction proved above.
By Lemma~\ref{lem:lp-validity} and Theorem~\ref{thm:fpt}, some
construction succeeds by $k=b(G)$. Stop at the first success
$\widehat k$, so $\widehat k\le b(G)$.
The preceding bounds give length at most
$(2+a+\varepsilon)b(G)$ in $n^{O_\varepsilon(1)}$ time.
For an arbitrary real $\varepsilon>0$, use any smaller positive rational
tolerance at most one.
Finally, $2+a=105607/50000=2.11214$; take
$\varepsilon=10^{-6}$ and use integrality of the returned length.
\end{proof}


\section{Hardness below two}\label{sec:hardness}

The reduction below produces connected bipartite graphs.

\begin{definition}
A Set Cover instance consists of a nonempty finite universe
$\mathcal U$ and an indexed family
$\mathcal S=(S_1,\ldots,S_M)$ of nonempty subsets of $\mathcal U$ whose
union is $\mathcal U$. Its optimum is
\[
 \tau(\mathcal U,\mathcal S)
 =\min\{|J|:J\subseteq[M],\ \bigcup_{j\in J}S_j=\mathcal U\}.
\]
\end{definition}

We use the following standard gap form of Set Cover hardness. It is the
only external hardness theorem used in the reduction.

\begin{lemma}[Set Cover gap {\cite[Corollary 1.5 and Section 6.2]{DS}}]
\label{lem:set-cover-gap}
For every fixed constant $\Gamma>1$, it is NP-hard to distinguish
\[
 \tau(\mathcal U,\mathcal S)\le k
 \quad\text{from}\quad
 \tau(\mathcal U,\mathcal S)>\Gamma k,
\]
where $1\le k\le M$.
\end{lemma}
The logarithmic gap supplied by the cited reduction implies every fixed
constant gap by restricting to sufficiently large instances. Instances
of bounded universe size can be solved by enumerating covered subsets;
thus discarding those instances does not remove the hardness. Empty
sets can be deleted without changing the optimum.

Fix an integer $\ell\ge3$ and an instance in
Lemma~\ref{lem:set-cover-gap} with
\begin{equation}\label{eq:hardness-constants}
 \Gamma=(\ell+1)(4\ell-2),\quad
 h=2k,\quad R=\ell h,\quad
 L_{\mathrm y}=R+h=2(\ell+1)k,\quad
 L_{\mathrm n}=2R-h=(4\ell-2)k.
\end{equation}
Here $L_{\mathrm y}$ and $L_{\mathrm n}$ will be the two graph-burning
thresholds.

\subsection{The graph and its coordinate labels}

Add a symbol $0$ not in $\mathcal U$. For each $e\in\mathcal U$, put
$I_e=\{j\in[M]:e\in S_j\}$, and put $I_0=[M]$. Each $I_e$ is nonempty.
For every $e\in\mathcal U\cup\{0\}$, introduce a new coordinate symbol
$\star_e$. All these coordinate symbols are distinct and are not in
$[M]$. Let $\mathbf e_\alpha$ denote the vector with one in coordinate
$\alpha$ and zero in the other coordinates.

For each $e\in\mathcal U\cup\{0\}$, form a disjoint copy of a graph
$K_e$ with vertex set
\begin{equation}\label{eq:simplex-vertices}
 V(K_e)=
 \left\{a\in\Z_{\ge0}^{\,I_e\cup\{\star_e\}}:
                        \sum_\alpha a_\alpha=\ell\right\}.
\end{equation}
Distinct vertices $a,b$ are adjacent exactly when
$b-a=\mathbf e_\alpha-\mathbf e_\beta$ for distinct coordinates
$\alpha,\beta$. Identify the vertices $\ell\mathbf e_j$ for the same
$j\in[M]$ across all graphs in which they occur; denote the resulting
vertex by $s_j$. Make no other identifications.
Vertices and edges of each $K_e$ are called coarse vertices and coarse
edges. Replace every coarse edge by a path of exactly $h$ edges, using new interior
vertices for each edge. The resulting graph is $G$.
Write $H_e$ for the image in $G$ of the subdivided $K_e$, and call
\[
 t_e=\ell\mathbf e_{\star_e}
\]
its terminal. The subgraph $H_0$ ensures connectedness; its terminal is
not an element of the original universe.

We use vectors with global coordinate set
\[
 \mathcal D=[M]\ \cup\ \{\star_e:e\in\mathcal U\cup\{0\}\}.
\]
At a coarse vertex $a$ of $K_e$ define $w(a)=ha$, extended by zero
outside its local coordinate set. In particular, $w(s_j)=R\mathbf e_j$
and $w(t_e)=R\mathbf e_{\star_e}$.
For an oriented coarse edge from $a$ to
$b=a+\mathbf e_\alpha-\mathbf e_\beta$, label its interior vertex at
distance $t\in[h-1]$ from $a$ by
\begin{equation}\label{eq:subdivision-label}
 w(x_t)=ha+t(\mathbf e_\alpha-\mathbf e_\beta).
\end{equation}
Changing the orientation gives the same label at the same vertex.
Labels at identified vertices agree. Every label is nonnegative and
has total coordinate sum $R$. Along each edge of $G$, every coordinate
changes in absolute value by at most one. Therefore
\begin{equation}\label{eq:coordinate-lower-bound}
 d_G(x,y)\ge |w_\alpha(x)-w_\alpha(y)|
 \quad(x,y\in V,\ \alpha\in\mathcal D).
\end{equation}

\begin{lemma}\label{lem:gadget-geometry}
The graph $G$ is connected and bipartite, and its size is polynomial
in the Set Cover input size for fixed $\ell$.
For every $e\in\mathcal U\cup\{0\}$ and $j\in I_e$,
\begin{equation}\label{eq:gadget-cover-radius}
 V(H_e)\subseteq B(s_j,R+h/2).
\end{equation}
\end{lemma}
\begin{proof}
A graph $K_e$ is connected because one can move its $\ell$ coordinate
units one at a time between coordinates. The graph $H_0$ contains every
$s_j$, and every other $H_e$ meets it at some $s_j$, proving
connectedness. Color every coarse vertex zero, and color an interior
vertex by the parity of its distance from an endpoint of its subdivided
edge. Since $h=2k$ is even, both endpoints have color zero and every edge
joins different colors. This is a bipartition.

The number of coarse vertices of $K_e$ is
$\binom{\ell+|I_e|}{\ell}\le(M+\ell+1)^\ell$, and each has degree at
most $(M+1)^2$. There are $|\mathcal U|+1$ such graphs and at most
$h-1\le2M$ new vertices per coarse edge. Enumeration of these vertices
and edges is polynomial-time for fixed $\ell$.

To prove~\eqref{eq:gadget-cover-radius}, move the $\ell-a_j$ units
outside coordinate $j$ of a coarse vertex $a$ to coordinate $j$.
This gives a path of $\ell-a_j$ coarse edges from $a$ to
$\ell\mathbf e_j$, hence distance at most $h(\ell-a_j)\le R$ in $G$.
If $x$ is at distance $t$ from one endpoint of a subdivided edge, its
distance from $s_j$ is at most
$\min\{R+t,R+h-t\}\le R+h/2$. This also covers all interior vertices.
\end{proof}

\subsection{Completeness and soundness}

\begin{lemma}\label{lem:hardness-completeness}
If $\tau(\mathcal U,\mathcal S)\le k$, then $b(G)\le L_{\mathrm y}$.
\end{lemma}
\begin{proof}
Let $J\subseteq[M]$ be a cover of size at most $k$.
Ignite its vertices $s_j$ at distinct rounds in $[k]$.
For every $e\in\mathcal U$ at least one such $j$ belongs to $I_e$.
For $e=0$, every chosen $j$ belongs to $I_0$, and $J$ is nonempty.
Lemma~\ref{lem:gadget-geometry} shows that the balls of radius $R+h/2$
at these sources cover $V(G)$. Since
$L_{\mathrm y}-k=R+h/2$, they fit their assigned rounds.
Lemma~\ref{lem:legalization} gives the asserted legal sequence.
\end{proof}

For a vertex $x$ that is not a shared vertex $s_j$, let
$\operatorname{own}(x)$ be its unique index
$e\in\mathcal U\cup\{0\}$ with $x\in V(H_e)$. For a shared vertex put
$\operatorname{own}(x)=\bot$, where $\bot$ is distinct from all indices.
This is well-defined because only the $s_j$ were identified.
Define
\begin{equation}\label{eq:heavy-coordinates}
 J(x)=\{j\in[M]:w_j(x)\ge h\}.
\end{equation}
Nonnegativity and the total sum $R=\ell h$ imply $|J(x)|\le\ell$.

\begin{lemma}\label{lem:ball-extraction}
For every $x\in V(G)$, the original elements $f\in\mathcal U$ whose
terminals lie in $B(x,L_{\mathrm n})$ can be covered by at most
$\ell+1$ sets of $\mathcal S$.
\end{lemma}
\begin{proof}
First consider $f\ne\operatorname{own}(x)$ with
$d_G(x,t_f)\le L_{\mathrm n}$. Any path from $x$ to $t_f$ meets a
shared vertex $s_j$ with $j\in I_f$, since only such vertices connect
$H_f$ to the rest of the graph. If $x$ is already a shared vertex of
$H_f$, use that vertex if necessary. Choose such a vertex along a
shortest path. The two portions of this path have lengths at least
$R-w_j(x)$ and $R$, respectively, by~\eqref{eq:coordinate-lower-bound}:
use coordinate $j$ for the first portion and $\star_f$ for the second.
Consequently
\[
 2R-w_j(x)\le d_G(x,t_f)\le 2R-h.
\]
Thus $j\in J(x)$ and $f\in S_j$.
All these elements are covered by the at most $\ell$ sets indexed by
$J(x)$. If $\operatorname{own}(x)\in\mathcal U$, add any one set
containing that element, which exists by the definition of the Set
Cover instance. This covers the only possible exceptional element.
\end{proof}

\begin{lemma}\label{lem:hardness-soundness}
If $\tau(\mathcal U,\mathcal S)>\Gamma k$, then
$b(G)>L_{\mathrm n}$.
\end{lemma}
\begin{proof}
A covering sequence of length at most $L_{\mathrm n}$ covers every
original terminal with at most $L_{\mathrm n}$ balls, each of radius
at most $L_{\mathrm n}$. Apply Lemma~\ref{lem:ball-extraction} to each
center. The extracted sets cover $\mathcal U$, and their number is at
most
\[
 (\ell+1)L_{\mathrm n}
 =(\ell+1)(4\ell-2)k=\Gamma k,
\]
contrary to the hypothesis.
\end{proof}

\begin{proof}[Proof of Theorem~\ref{thm:hardness}]
Use Lemma~\ref{lem:set-cover-gap} with the constant in
\eqref{eq:hardness-constants} and construct $G$ as above.
Lemmas~\ref{lem:gadget-geometry},
\ref{lem:hardness-completeness}, and
\ref{lem:hardness-soundness} prove the claimed polynomial-time
reduction on connected bipartite graphs.
For a fixed $\alpha\in[1,2)$, choose a fixed $\ell\ge3$ such that
\[
 \alpha\le\frac{L_{\mathrm n}}{L_{\mathrm y}}
          =2-\frac3{\ell+1}.
\]
An $\alpha$-approximation returns length at most $L_{\mathrm n}$ in
the positive case, whereas every legal sequence has length greater
than $L_{\mathrm n}$ in the negative case. It would decide this
NP-hard promise problem in polynomial time.
\end{proof}


\appendix
\section{Exact verification of the numerical inequality}
\label{app:certificate}

This appendix specifies finite rational bounds proving
Lemma~\ref{lem:numerical}. The integer programs in
Sections~\ref{app:python} and~\ref{app:cpp} evaluate these bounds.
All exponential functions in the verification have arguments in
$[0,W]\subset[0,2]$.

\subsection{Outward rational evaluation of the exponential functions}

Set
\[
 \mathsf S=10^{12},\qquad \mathsf T=10^{40},\qquad
 N_r=N_t=65536,\qquad N_q=524288,\qquad C_*=100000.
\]
For $d\in\Z_{\ge0}$ define the rational polynomials
\[
 E_d(z)=\sum_{i=0}^d\frac{(-1)^iz^i}{i!},
 \qquad V_d(z)=\sum_{i=0}^d\frac{(-1)^iz^i}{(i+1)!}.
\]
Taylor's theorem with its signed remainder gives, for $z\ge0$,
\begin{equation}\label{eq:taylor-bounds}
 E_{21}(z)\le e^{-z}\le E_{20}(z),\qquad
 V_{21}(z)\le\phi(z)\le V_{20}(z).
\end{equation}
The second pair also follows by integrating the first pair at $tz$
over $t\in[0,1]$.

For an upper bound, round every coefficient of $E_{20}$ or $V_{20}$
up to a multiple of $\mathsf T^{-1}$, round the argument $z$ down
to a multiple of $\mathsf T^{-1}$, and evaluate by Horner's rule,
rounding every multiplication up to a multiple of $\mathsf T^{-1}$.
Finally round up to a multiple of $\mathsf S^{-1}$ and take the minimum
with one. Denote the resulting bounds by $U_E(z)$ and $U_\phi(z)$.
For a lower exponential bound use $E_{21}$, round coefficients down,
round the argument up, round Horner multiplications down, finally
round down to a multiple of $\mathsf S^{-1}$, and take the maximum
with zero. Denote it by $L_E(z)$.

These argument-rounding directions are valid because the rounded
polynomials are decreasing on $[0,2]$. Here are explicit bounds proving
that fact. The elementary estimate $e^2<8$ gives $e^{-2}>1/8$, and
\[
 E_{19}(z)\ge e^{-z}-2^{20}/20!>1/9\quad(0\le z\le2).
\]
Thus $E_{20}'=-E_{19}<-1/9$, $E_{21}'=-E_{20}<-1/8$, and
$V_{20}'(z)=-\int_0^1tE_{19}(tz)\,dt<-1/18$.
Coefficient rounding changes any of these derivatives in absolute
value by at most
$\mathsf T^{-1}\sum_{i=1}^{21}i2^{i-1}<10^{-31}$.
Indeed,
\[
 e\le1+1+\tfrac12+\tfrac16+
       \tfrac1{24}\sum_{t\ge0}5^{-t}=\tfrac{87}{32}<\tfrac{11}{4},
\]
so $e^2<8$.
At a fixed nonnegative argument, directed Horner rounding preserves
upper and lower bounds at every step. Consequently
\begin{equation}\label{eq:oracle-bounds}
 0\le L_E(z)\le e^{-z}\le U_E(z)\le1,\qquad
 \phi(z)\le U_\phi(z)\le1
 \quad(0\le z\le2).
\end{equation}
The programs implement ceilings and floors of signed integers, so
these assertions include negative intermediate Horner values.

\subsection{Primitive bounds and enclosure of every cut}

Use the primitive grid $x_j=j/N_q$ for $0\le j\le3N_q$.
Evaluate $g(x_j)$ exactly by~\eqref{eq:g-finite} and define
\[
 q_j^-=L_E(g(x_j)),\qquad q_j^+=U_E(g(x_j)).
\]
On an open cell $(x_j,x_{j+1})$, the constants $q_j^-$ and
$q_{j+1}^+$ bound $q_P$ from below and above, respectively.
Let $P^-(x)$ and $P^+(x)$ be the integrals from zero to $x$ of
these lower and upper step functions. At grid points, use $q_j^-$
and $q_j^+$ for point evaluations. Integration at a partial cell is
linear interpolation of the appropriate step primitive.
For $0\le u\le v\le3$ these give
\begin{equation}\label{eq:primitive-bounds}
 P^-(v)-P^-(u)\le\int_u^vq_P(x)\,dx
                    \le P^+(v)-P^+(u).
\end{equation}
The programs round partial evaluations outward at denominator
$\mathsf H=\mathsf S N_q$: a subtracted starting primitive is rounded
in the opposite direction from the ending primitive.

For $r\in[0,1]$ and $e\ge0$, set
\[
 Q_r(e)=\int_{x<e}q_P(x)\,d\nu_r(x),\qquad
 T_r(e)=\nu_r([e,+\infty)),\qquad
 H_r(e)=Q_r(e)+\phi(g(e))T_r(e).
\]
In particular,
\begin{equation}\label{eq:exact-tail}
 T_r(e)=W-A_0\ind_{e>r}
       -\sum_{v=1}^3p_v[e-r-c_{v-1}]_0^{c_v-c_{v-1}}.
\end{equation}
Compute lower and upper bounds $Q_r^-(e),Q_r^+(e)$ by the atom
$A_0q_P(r)\ind_{r<e}$ and by applying~\eqref{eq:primitive-bounds}
on each of the three shifted density intervals. A density interval
whose lower endpoint is at least $e$ contributes zero.
All additions and multiplications by nonnegative densities are
rounded outward. Then
\begin{equation}\label{eq:Q-enclosure}
 Q_r^-(e)\le Q_r(e)\le Q_r^+(e).
\end{equation}
For example, the exact unrounded continuous part is
\[
 \sum_{\substack{v\in[3]\\r+c_{v-1}<e}}
       p_v\int_{r+c_{v-1}}^{\min\{e,r+c_v\}}q_P(x)\,dx.
\]
The atom is handled separately.

Use radius parameters $\rho_i=i/N_r$, $i\in[N_r]$, and cut points
$\theta_j=j/N_t$, $j\in\{0,\ldots,3N_t\}$.
Put $J_c=\lceil c_3N_t\rceil$.
Because $N_t$ is divisible by $N_r$, the only atom of
$\nu_{\rho_i}$ is always a cut point.
For a grid cut, an upper bound for $H_r(\theta_j)$ is
\begin{equation}\label{eq:cut-point}
 Q_r^+(\theta_j)+U_\phi(g(\theta_j))T_r(\theta_j).
\end{equation}
For every cut $e\in(\theta_j,\theta_{j+1})$, monotonicity gives the
upper bound
\begin{equation}\label{eq:cut-cell}
 Q_r^+(\theta_{j+1})
 +U_\phi(g(\theta_{j+1}))
          \bigl(T_r(\theta_j)-A_0\ind_{\theta_j=r}\bigr).
\end{equation}
The subtraction removes the atom at the left endpoint: it has moved
from the tail into the strict prefix on the open cell.
For $e\ge c_3$, the value of $H_r(e)$ is constant, since both
$q_P(x)$ and $\phi(g(x))$ equal one for $x\ge c_3$.

By the definition of $\Phi$,
\[
 K(r,\theta)=\sup_{e\ge\theta}H_r(e)-Q_r(\theta).
\]
Take a reverse running maximum of~\eqref{eq:cut-point} and
\eqref{eq:cut-cell}, including an upper bound for
$H_r(c_3)=Q_r(c_3)+T_r(c_3)$. Subtract $Q_r^-(\theta_j)$.
Taking the minimum of this upper bound with $T_r(\theta_j)$,
and taking its maximum with zero, still gives an upper bound for
$K(r,\theta_j)$. If $\theta_j\ge c_3$, use the exact value
$K(r,\theta_j)=T_r(\theta_j)$.
Denote the resulting outward-rounded bound by $\overline K_{i,j}$.

Monotonicity in the radius parameter, proved in
Lemma~\ref{lem:kernel-properties}, yields
\begin{equation}\label{eq:MN-certificate}
 M_N(\theta_j)\le\frac1{N_r}\sum_{i=1}^{N_r}\overline K_{i,j}.
\end{equation}
The computation uses this bound for $j\le J_c$. For larger $j$ the
other bound below suffices.

\subsection{The aggregate bound and the capacity minimum}

An exact cumulative integral for $\Lambda$, valid for $x\ge0$, is
\begin{equation}\label{eq:aggregate-cdf}
\begin{split}
 C_\Lambda(x)={}&A_0[x]_0^1\\
 &+\frac12\sum_{v=1}^3p_v\left(
 (x-c_{v-1})_+^2-(x-c_v)_+^2
 -(x-1-c_{v-1})_+^2+(x-1-c_v)_+^2\right).
\end{split}
\end{equation}
It follows by integrating~\eqref{eq:aggregate-density}.
Thus the masses of all cut cells are exact rationals.
Since $e^{-D(x)}$ is nondecreasing and $\Lambda$ is supported below
three,
\begin{equation}\label{eq:MO-certificate}
 M_O(\theta_j)\le
 \sum_{v=j}^{3N_t-1}
 U_E(D(\theta_{v+1}))
       \bigl(C_\Lambda(\theta_{v+1})-C_\Lambda(\theta_v)\bigr).
\end{equation}
The programs round each cell mass and each product upward at
denominator $\mathsf H$.

Let $\overline M_j$ be the smaller of the outward-rounded bounds
\eqref{eq:MN-certificate} and~\eqref{eq:MO-certificate}, using
\eqref{eq:MO-certificate} alone when $j>J_c$.
Then $M(\theta_j)\le\overline M_j$.
Define
\[
 C(\theta)=2+a-\theta-M(\theta),\qquad
 c_j^-=2+a-\theta_j-\overline M_j.
\]
The letter $C$ here denotes a real-valued function, not a graph
component. It is used only in this verification.
For a cut cell $[\theta_{j-1},\theta_j]$, the rational number
\begin{equation}\label{eq:density-cell-upper}
 \beta_j=A_0\ind_{\theta_{j-1}<1}
  +\sum_{v=1}^3p_v
    \pos{\min\{1,\theta_j-c_{v-1}\}
                   -\max\{0,\theta_{j-1}-c_v\}}
\end{equation}
is an upper bound for the density of $\Lambda$ on its interior.
If $\beta_j\le1$, equation~\eqref{eq:envelope-increment} proves
that $C$ is nonincreasing on that cell. Its minimum there is at the
right endpoint and is at least $c_j^-$.
In every other cell, monotonicity of $M$ gives the lower bound
$c_{j-1}^- - 1/N_t$ throughout the cell.

Start with $b_0=c_0^-$ and recursively define
\begin{equation}\label{eq:capacity-recursion}
 b_j=\min\left\{b_{j-1},c_j^-,
 \begin{cases}
 c_j^-,&\beta_j\le1,\\
 c_{j-1}^- -1/N_t,&\beta_j>1
 \end{cases}\right\}\quad(1\le j\le2N_t).
\end{equation}
This is a lower bound for the minimum of $C$ on $[0,\theta_j]$,
which equals $A(2-\theta_j)$.
The entire recursion is on the exact scale $\mathsf H$.

Changing variables to $\theta=2-s$ in~\eqref{eq:integral-J} gives
a nondecreasing integrand:
\[
 J=\frac12\int_0^2\frac{e^{-g(\theta/2)}}{A(2-\theta)}\,d\theta.
\]
Consequently, if all $b_j>0$, an upper bound is
\begin{equation}\label{eq:J-finite-certificate}
 J\le\frac1{2N_t}\sum_{j=1}^{2N_t}
           \frac{U_E(g(\theta_j/2))}{b_j}.
\end{equation}
Each summand is rounded upward to a multiple of $\mathsf S^{-1}$.
The resulting sum has numerator $131067519775194526$ and denominator
$2N_t\mathsf S=131072000000000000$.
The minimum of the computed $b_j$ is
$54337893542773910/\mathsf H$.
Since $A$ is nondecreasing, these grid bounds also imply its uniform
lower bound on $[0,2]$. This completes the proof of
Lemma~\ref{lem:numerical}.

\subsection{Reproduction and integer ranges}

Compiling this LaTeX file writes three files in the compilation
directory:
\begin{itemize}[nosep]
\item \texttt{graph\_burning\_profile.json};
\item \texttt{graph\_burning\_certificate.py};
\item \texttt{graph\_burning\_certificate.cpp}.
\end{itemize}
They are also displayed below.
With Python 3 and a GNU C++17 compiler supporting signed 128-bit
integers, run
\begin{quote}\ttfamily\small
python3 graph\_burning\_certificate.py graph\_burning\_profile.json
\end{quote}
on a little-endian machine with eight-byte signed integer arrays.
The profile, tables, and decisions use integer or rational arithmetic.
Elapsed-time measurements are only printed.

The Python program has arbitrary-precision integers.
In the C++ program put $\mathsf H=\mathsf S N_q$ as above.
Every stored signed 64-bit arithmetic value, apart from its initial
sentinel, has absolute value less than
$6\mathsf H<2^{63}$. To check this, primitive sums are at most
$3\mathsf H$; the scaled values representing $Q_r(e)$ and $T_r(e)$ are each less than
$2\mathsf H$ including outward rounding; raw capacity uses at most
$3\mathsf H$, $2\mathsf H$, and $2\mathsf H$ in its three terms.
The other stored quantities are smaller.
The largest signed 128-bit products occur in~\eqref{eq:aggregate-cdf}.
Its three density contributions and its atom contribution are
nonnegative, and their sum is at most
\[
 2(W\mathsf S)N_q(C_*N_t)^2<2^{126}.
\]
The separate squared positive parts are at most $9(C_*N_t)^2$.
Products for integration in the radius parameter are less than $10^{31}$,
and the accumulated sums over that parameter are less than $2N_r\mathsf H$.
All are within signed 128-bit range. The Python program checks the
displayed $2^{126}$ bound before running the C++ program.
The C++ program checks the exact total aggregate mass, positive
capacities, and the final rational inequalities.

\subsection{The rational profile}\label{app:profile}

\lstinputlisting[language={},caption={The rational enlargement profile.}]{graph_burning_profile.json}

\subsection{Table generation}\label{app:python}

\lstinputlisting[language=Python,caption={Directed rational evaluation and table generation.}]{graph_burning_certificate.py}

\subsection{Finite-sum evaluation}\label{app:cpp}

\lstinputlisting[language=C++,caption={Exact bounds for the cut supremum, capacity minimum, and integral.}]{graph_burning_certificate.cpp}

\begin{thebibliography}{9}

\bibitem{BNZ}
J.~Blauth, C.~N\"obel, and R.~Zenklusen.
\newblock Toward optimal approximations for resource-minimization for
fire containment on trees and non-uniform $k$-center.
\newblock \emph{arXiv:2512.17049}, 2025.
\newblock \url{https://arxiv.org/abs/2512.17049}.

\bibitem{DS}
I.~Dinur and D.~Steurer.
\newblock Analytical approach to parallel repetition.
\newblock In \emph{Proceedings of the 46th Annual ACM Symposium on
Theory of Computing}, pages 624--633, 2014.
\newblock Full version: \emph{arXiv:1305.1979v3}.
\newblock \url{https://doi.org/10.1145/2591796.2591884}.

\bibitem{Khachiyan}
L.~G.~Khachiyan.
\newblock Polynomial algorithms in linear programming.
\newblock \emph{USSR Computational Mathematics and Mathematical Physics},
20(1):53--72, 1980.
\newblock \url{https://doi.org/10.1016/0041-5553(80)90061-0}.

\end{thebibliography}
\end{document}

